\documentclass[pdflatex,sn-mathphys-ay]{sn-jnl}

\usepackage[utf8]{inputenc}%
\usepackage[T1,T2A]{fontenc}%
\usepackage[english,russian]{babel}%
\usepackage{graphicx}%
\usepackage{multirow}%
\usepackage{amsmath,amssymb,amsfonts}%
\usepackage{amsthm}%
\usepackage{mathrsfs}%
\usepackage[title]{appendix}%
\usepackage{xcolor}%
\usepackage{textcomp}%
\usepackage{manyfoot}%
\usepackage{booktabs}%
\usepackage{array}%
\usepackage{longtable}%
\usepackage{calc}%
\usepackage{newunicodechar}%
\providecommand{\tightlist}{\setlength{\itemsep}{0pt}\setlength{\parskip}{0pt}}
\newunicodechar{§}{\S}
\newunicodechar{→}{\ensuremath{\rightarrow}}
\newunicodechar{↔}{\ensuremath{\leftrightarrow}}
\newunicodechar{×}{\ensuremath{\times}}
\newunicodechar{·}{\ensuremath{\cdot}}
\newunicodechar{⊃}{\ensuremath{\supset}}
\newunicodechar{°}{\ensuremath{^{\circ}}}
\newunicodechar{≠}{\ensuremath{\neq}}
\newunicodechar{≥}{\ensuremath{\geq}}
\newunicodechar{≤}{\ensuremath{\leq}}
\newunicodechar{≈}{\ensuremath{\approx}}
\newunicodechar{∼}{\ensuremath{\sim}}
\newunicodechar{±}{\ensuremath{\pm}}
\newunicodechar{∓}{\ensuremath{\mp}}
\newunicodechar{∈}{\ensuremath{\in}}
\newunicodechar{∉}{\ensuremath{\notin}}
\newunicodechar{⊥}{\ensuremath{\perp}}
\newunicodechar{∥}{\ensuremath{\parallel}}
\newunicodechar{∧}{\ensuremath{\wedge}}
\newunicodechar{∨}{\ensuremath{\vee}}
\newunicodechar{¬}{\ensuremath{\neg}}
\newunicodechar{≫}{\ensuremath{\gg}}
\newunicodechar{≪}{\ensuremath{\ll}}
\newunicodechar{⇒}{\ensuremath{\Rightarrow}}
\newunicodechar{⇔}{\ensuremath{\Leftrightarrow}}
\newunicodechar{∞}{\ensuremath{\infty}}
\newunicodechar{←}{\ensuremath{\leftarrow}}
\newunicodechar{∂}{\ensuremath{\partial}}
\newunicodechar{∇}{\ensuremath{\nabla}}
\newunicodechar{√}{\ensuremath{\surd}}
\newunicodechar{∝}{\ensuremath{\propto}}
\newunicodechar{≅}{\ensuremath{\cong}}
\newunicodechar{≡}{\ensuremath{\equiv}}
\newunicodechar{÷}{\ensuremath{\div}}
\newunicodechar{∅}{\ensuremath{\emptyset}}
\newunicodechar{¹}{\textsuperscript{1}}
\newunicodechar{²}{\textsuperscript{2}}
\newunicodechar{³}{\textsuperscript{3}}
\newunicodechar{⁴}{\textsuperscript{4}}
\newunicodechar{⁵}{\textsuperscript{5}}
\newunicodechar{⁶}{\textsuperscript{6}}
\newunicodechar{⁷}{\textsuperscript{7}}
\newunicodechar{⁸}{\textsuperscript{8}}
\newunicodechar{⁹}{\textsuperscript{9}}
\newunicodechar{⁰}{\textsuperscript{0}}
\newunicodechar{«}{\guillemotleft}
\newunicodechar{»}{\guillemotright}

\raggedbottom

\begin{document}

\title[Витальность как самооплаченное самомоделирование]{Витальность как эффективность самооплаченного самомоделирования}

\author*[1]{\fnm{Александр} \sur{Андриишин}}\email{alexander@andriishin.com}

\affil*[1]{\orgname{Independent Researcher}, \orgaddress{\city{Москва}, \country{Россия}}}

\abstract{Сложность вездесуща, живое --- редкость: звёзды, ураганы, кристаллы и языковые модели богато структурированы и далеки от равновесия, но ни одна из этих систем не оплачивает собственной диссипацией удерживаемую ею модель мира. Замкнутая \emph{самооплата} отделяет витальное от просто сложного; здесь предложена её количественная операционализация. Витальность определяется парой условий. Первое --- предсказательная эффективность самомоделирования: доля удерживаемой системой памяти о среде, которая ещё предсказывает её будущее. Второе --- самооплата: термодинамическая цена удержания устаревшей памяти обязывает систему к диссипации, а искажение модели возвращается её распадом. Требование, чтобы модель и оплачивающая её диссипация принадлежали одной физической границе, отличает шкалу от известных программ --- assembly theory, интегрированной информации, телеодинамики, принципа свободной энергии: каждая берёт одну ось витальности всерьёз, но ни одна не требует этого тождества границ. Двухступенчатая процедура --- структурный скрининг по триаде инвариантов и витальная верификация --- применена к шести камертонам от звезды до мегаполиса, с биосферой как контрольным пределом и воспроизводимым расчётом на бактериальном хемотаксисе (\emph{Escherichia coli}). Картина асимметрична: структурный потенциал почти универсален, замкнутые петли самооплаты редки. Отсюда следствия: поиск внеземной жизни требует обнаружения таких петель поверх химических биосигнатур; этический статус искусственных систем --- вопрос о появлении собственной платы за модель; граница между сложным и витальным проводится по положительной эффективности при замкнутой петле. Теория развёрнута для стационарного режима; нестационарное обобщение дано в сопровождающей работе (Andriishin 2026). Программа открыта к опровержению по нескольким независимым линиям, включая предсказанную степенную зависимость темпа эволюционной адаптации от эффективности самомоделирования.}

\keywords{определение жизни, витальность, самооплата, ностальгия, predictive information, термодинамика вычислений}

\maketitle

\section{Введение}\label{ux432ux432ux435ux434ux435ux43dux438ux435}

Концепт «живой организм» в современной биологии не имеет общепринятого формального определения. Исторически он опирался на клеточный субстрат --- углеродную химию, мембранную организацию, метаболическое замыкание, --- но его границы размыты как снизу, так и сверху. Снизу: вирусы лишены автономного метаболизма, но участвуют в эволюционной динамике (Pradeu 2018); прионы и репликаторные РНК-системы воспроизводятся без клеточной инфраструктуры; протоклеточные и абиогенетические переходные формы не позволяют провести резкую границу между химией и биологией (Plante 2025). Сверху: биосферы, симбиотические сообщества и антропогенные мегаструктуры демонстрируют признаки коллективной авто-регуляции и эволюционной адаптации без единого клеточного референта (Mossio et al.~2016). В ответ на эти трудности философия биологии последних двух десятилетий выработала устойчивый \textbf{эпистемологический плюрализм}:

\begin{itemize}
\tightlist
\item
  реестры из ста и более конкурирующих определений жизни без признаков сходимости (Trifonov 2011);
\item
  аргумент о необходимости «теории жизни» вместо словарного определения (Cleland 2019);
\item
  анти-дефиниционистская линия «жизнь-как-историческая-линия»: жизнь как единая монофилетическая ветвь от LUCA, а не список признаков (Mariscal and Doolittle 2020);
\item
  концепция организма как биологически индивидуированной системы (Bich and Green 2018; Pradeu 2018).
\end{itemize}

Ни одно определение не получает единодушного признания; нарастает запрос на операциональные критерии, способные приложиться к нестандартным кандидатам --- астробиологическим находкам, искусственным системам сложной обработки информации, традиционно называемым когнитивными, гибридным экологическим комплексам (Chirumbolo and Vella 2026).

Настоящая работа предлагает на эту роль операциональную меру --- безразмерную предсказательную эффективность самомоделирования \(\eta_v = I_{\text{pred}}/I_{\text{mem}}\), а вердикт витальности --- как пару \((S, \eta_v)\) из предиката самооплаты \(S\) и положительной эффективности \(\eta_v > 0\). Прежде чем развернуть эту конструкцию (§ 2), очертим ландшафт программ, на фоне которого она строится.

Витальность формализует линия технических программ, каждая берущая одну её сторону: термодинамическую («жизнь питается негэнтропией», Schrödinger 1944), байесовскую (FEP: Friston 2005, 2010), комбинаторную сложность сборки (assembly theory: Marshall et al.~2017; Sharma et al.~2023), каузальную замкнутость информации (\(\Phi\): Tononi 2004; Tononi et al.~2016), самовозникновение норм (телеодинамика: Deacon 2011), масштаб целеподобного поведения (cognitive light cone: Levin 2019). Это не взаимоисключающие гипотезы: каждая теряет дискриминационную силу на границах --- FEP применим к бытовому термостату (Baltieri and Buckley 2019) и регулятору Уатта (Bruineberg et al. 2022), узкое NASA-определение (Joyce 1994; Cleland and Chyba 2002) молчит о субстратной общности (детальное сопоставление --- § 4). Отсюда мост к нашей мере: соединить термодинамическую плату и качество удерживаемой модели в одну безразмерную шкалу.

Параллельная программа \textbf{происхождения жизни} --- гиперцикл (Eigen 1971; Eigen and Schuster 1979), autocatalytic sets и RAF-сети (Hordijk and Steel 2004; Hordijk 2019), dissipative adaptation (Perunov, Marsland, and England 2016), major transitions (Maynard Smith and Szathmáry 1995), геохимические сценарии (Pross 2012; Smith and Morowitz 2016; Kauffman 2019) --- операционализирует переход abiotic→biotic как самоподдерживающуюся информационно-энергетическую петлю над геохимическим потоком; \(\eta_v\) делает то же в термодинамическом регистре, где положительность \(\eta_v\) --- необходимое условие момента origin of life (§ 4.8; § 6).

Те же пределы --- в астробиологии: биосигнатуры по неравновесности (\(\text{O}_2\)/\(\text{CH}_4\)), изотопам (\(\delta^{13}\text{C}\)), хиральности имеют абиотические альтернативы, и дебаты об arsenate-claim штамма GFAJ-1 (\emph{Halomonadaceae}) (Wolfe-Simon et al.~2011 {[}retracted by \emph{Science} 2025{]}; Reaves et al.~2012), метане Марса (Webster et al.~2018; Korablev et al.~2019), фосфине Венеры (Greaves et al.~2021; Villanueva et al.~2021) без витального критерия сводятся к интерпретации частных прокси --- как и спор о витальности языковых моделей.

С другой стороны к тому же вопросу подходит AI alignment: при каких архитектурных условиях оптимизирующий агент начинает оптимизировать собственное продолжение (Bostrom 2014; Hubinger et al.~2019; Carlsmith 2022). Объект пересекающийся, углы разные: AI-safety описывает поведенческие признаки петли, а \(\eta_v\) и предикат \(S\) --- её термодинамическую ко-характеристику (структурная параллель, не каузально-генеративная связь; § 3.4).

Настоящая работа определяет витальность как эффективность удержания собственной модели собственной диссипацией.

\emph{Самомоделирование} здесь используется в смысле, отличном от феноменальной self-model Метцингера (Metzinger 2003): оно обозначает структурное условие, при котором модель среды, оплачивающая её диссипация и носитель, страдающий от её ошибок, принадлежат одной физической системе --- ``self'' есть носитель модели и её цены, не содержание модели. Аналогично, \emph{витальность} здесь --- строго операциональное информационно-термодинамическое свойство (предсказательная эффективность самомоделирования \(\eta_v\)), не \emph{vis vitalis} виталистической традиции.

Центральный объект --- безразмерное отношение

\[\eta_v = \frac{I_{\text{pred}}}{I_{\text{mem}}} \in [0, 1],\]

где \(I_{\text{pred}}\) --- predictive information о будущем собственной среды, удерживаемая внутренней структурой системы, а \(I_{\text{mem}}\) --- полная информация той же структуры о прошлой и текущей траектории среды (точные определения и связь с классической формулировкой Биалека--Неменмана--Тишби (Bialek et al.~2001) --- § 2.1). Отношение \(\eta_v\) есть доля удерживаемой памяти, которая ещё предсказательна; дополнительная доля \(\nu = 1 - \eta_v\) --- \emph{информационная ностальгия}, балласт устаревшей памяти. Граница \(\eta_v \in [0, 1]\) для \textbf{пассивного} самомоделирования --- следствие неравенства об обработке данных (DPI), не энергетическая нормировка и не соглашение о масштабе (активный случай --- открытый пункт, § 2.1). Термодинамика входит \textbf{отдельным якорем}: по термодинамике предсказания (Still et al.~2012, § 2.2) удержание непредсказательной памяти обязывает систему к диссипации не ниже \(k_B T\) на каждый такой бит --- здесь и только здесь появляется ландауэровский масштаб \(k_B T \ln 2\), как коэффициент строгой границы, а не курс обмена энергии на биты. Витальность есть пара условий: положительная эффективность \(\eta_v > 0\) \textbf{при} замкнутой петле самооплаты (§ 2.7).

Программа разворачивается по трём конструктивным шагам. Первый --- операциональное определение витальной шкалы как \(\eta_v = I_{\text{pred}}/I_{\text{mem}}\) (свежесть удерживаемой модели), явно отличаемой и от индексов сложности (assembly index, \(\Phi\), \(I_v\) как композит трёх ёмкостей), и от предиката самооплаты \(S\): положительный \(\eta_v\) возможен и у системы без собственной платы (как у языковой модели), но витальность требует конъюнкции \(\eta_v > 0\) и \(S = 1\). Второй --- разделение двух методологических уровней: структурного скрининга через композит \(I_v = \sqrt[3]{T_v C_v S_v}\) (термодинамическая \(T_v\), информационная \(C_v\), топологическая \(S_v\) ёмкости; § 2.3) и витальной верификации через \(\eta_v\) (вычислительная экономия, не логическое усиление). Третий --- явная конвенциональность границы системы как методологический инструмент: \(\eta_v\) зависит от выбора границы, и эта зависимость продуктивна --- она позволяет различить «витальность LLM как агента» (равна нулю) и «витальность корпорации» (отлична от нуля, но относится к другому объекту).

Новизна --- не в отдельных компонентах. Predictive information (Bialek et al.~2001) и ландауэровский предел (Landauer 1961) известны более полувека. Самомоделирование восходит к anticipatory systems Розена (Rosen 1985); требование собственной платы за модель сходится с skin-in-the-game Талеба (Taleb 2018) и autopoiesis Матураны--Варелы (Maturana and Varela 1980). Новизна --- в архитектуре их соединения по трём пунктам. Специфическое отношение \(I_{\text{pred}}/I_{\text{mem}}\), насколько нам известно, не использовалось как мера витальности в линии Биалека--Неменмана--Тишби. Требование, что модель и оплачивающая её диссипация принадлежат одной физической границе (self-payment), с замкнутой петлёй возврата ошибки, в явной форме как критерий витальности, насколько нам известно, ранее не формализовано --- и Ландауэр, и Биалек работали без него. Интерпретация полученного отношения как витальной шкалы --- самостоятельный шаг, который ни одна из исходных рамок не делала.

Ближайшую современную линию --- термодинамику предсказания и обучения --- требуется отстроить отдельно, поскольку она оперирует родственными величинами. Голдт и Зайферт (Goldt and Seifert 2017) формализуют термодинамику обучения правилу в нейронной сети, вводя эффективность обучения как отношение приобретённой информации к её термодинамической стоимости, ограниченное единицей; Такахаси и Хаяси (Takahashi and Hayashi 2026) развивают КПД приобретения структурной информации о среде (\(\eta_E\), бит-на-джоуль) и требование граничного замыкания --- оплату низкоэнтропийного ресурса, пересекающего границу системы. Обе работы делят с настоящей общий фундамент (Landauer 1961; Still et al.~2012; Goldt and Seifert 2017), поэтому отстройка ведётся не по основанию, а по предмету, размерности и интерпретации, и линии комплементарны, а не конкурентны. Размерность: их \(\eta_E\) --- бит-на-джоуль о \emph{среде}, эффективность приобретения внешней структуры; \(\eta_v = I_{\text{pred}}/I_{\text{mem}}\) безразмерна, то есть: доля удерживаемой \emph{самомодели}, ещё остающаяся предсказательной, --- мера удержания и устаревания (ностальгии), не приобретения. Статус самооплаты: их граничное замыкание есть соглашение об учёте при выборе границы, тогда как предикат \(S\) (со-локация модели, оплачивающей её диссипации и петли возврата ошибки; § 2.7) --- конститутивный критерий витальности с независимой термодинамической опорой на бонд Still. Наконец, граница \(\eta_v \in [0, 1]\) следует из неравенства об обработке данных, не из стохастико-термодинамического баланса бит-на-джоуль, а витальной интерпретации пары \((S, \eta_v)\) термодинамика обучения не предлагает.

Предлагаемое определение не есть метафизическая декларация. \(\eta_v\) --- операциональное представление витальности как режима существования, не его определение по существу. Стандартное философское различение между референтом теории (свойство системы) и операционализацией (отношение между системой и измерительной рамкой) аналогично различению «температура как мера тепловой энергии» и «логарифм расширения ртути на 100 делений», или «сила землетрясения» и «шкала Рихтера». \(\eta_v\) --- операциональная мера, дисциплинированная конкретным определением.

Настоящая работа сознательно ограничивается \textbf{стационарным режимом}, где допущение конечности памяти и эргодичности среды стабилизирует объём удерживаемой памяти и делает информационную ностальгию стационарной. В этом режиме формализм даёт резкие демаркационные ответы и проверяемые предсказания. Нестационарный режим --- растущая память, дрейфующая среда, кумулятивный бюджет --- выходит за рамки данной работы и развит в сопровождающей работе (Andriishin 2026). Там доказывается Лемма 2 о неизбежности информационной ностальгии в OU-классе (Орнштейна--Уленбека) медленного дрейфа при перечисленных там технических условиях. Формулируется также гипотеза об адиабатической асимптотике \(\dot{\eta_v}^{\text{excess}}\) для систем без периодического сброса памяти. Применимость аппарата данной работы ограничена системами с эргодической средой и конечной памятью; для биосфер на эоновых окнах, обучающихся ИИ-систем под дрейфом и цивилизаций с растущим архивом --- обращаться к нестационарному расширению.

Структура статьи: § 2 разворачивает формальную рамку; § 3 применяет процедуру к шести камертонным системам (диагностическим тест-кейсам, каждый изолирует одно несущее условие); § 4 сопоставляет с соседними программами (assembly theory, IIT, телеодинамика, FEP) и позиционирует относительно линий организационного замыкания, происхождения жизни и современных демаркационных дебатов; § 5 формулирует условия фальсификации; § 6 обсуждает следствия для астробиологии, ИИ-этики, сложных систем; § 7 подводит итог.

\section{Теоретическая рамка}\label{ux442ux435ux43eux440ux435ux442ux438ux447ux435ux441ux43aux430ux44f-ux440ux430ux43cux43aux430}

\subsection{Центральная величина: предсказательная эффективность самомоделирования}\label{ux446ux435ux43dux442ux440ux430ux43bux44cux43dux430ux44f-ux432ux435ux43bux438ux447ux438ux43dux430-ux43fux440ux435ux434ux441ux43aux430ux437ux430ux442ux435ux43bux44cux43dux430ux44f-ux44dux444ux444ux435ux43aux442ux438ux432ux43dux43eux441ux442ux44c-ux441ux430ux43cux43eux43cux43eux434ux435ux43bux438ux440ux43eux432ux430ux43dux438ux44f}

Живая система удерживает во внутренней структуре неявную модель своей среды. Часть запомненного ещё предсказывает будущее среды --- она полезна; остальное есть след прошлого, от которого среда уже ушла, --- балласт. Витальность измеряется тем, \textbf{какая доля удерживаемой памяти остаётся предсказательной}: чем выше доля, тем «живее» петля самомоделирования, а чистый балласт без предсказания --- уже не работающая модель.

Формализация этой интуиции опирается на две информационные величины. Пусть \(M_t\) --- конфигурация тех внутренних степеней свободы системы \(X\) в момент \(t\), которые проходят counterfactual ablation (§ 2.2) как принадлежащие моделирующему контуру (функциональный канал, не полная внутренняя конфигурация --- для бактерии это паттерн метилирования Tar-рецепторов, а не все \(\sim 10^{10}\) степеней свободы клетки), а \(X_E\) --- сигнал её среды. В духе Still, Sivak, Bell, Crooks (2012) вводятся

\begin{equation}I_{\text{mem}}(t) := I\!\left(M_t;\; X_E^{\le t}\right), \qquad
  I_{\text{pred}}(t,\tau) := I\!\left(M_t;\; X_E^{[t,\,t+\tau]}\right). \tag{1a}\end{equation}

Здесь \(I_{\text{mem}}\) --- \textbf{память}: сколько внутреннее состояние знает о прошлой и текущей траектории среды (и эту память правомерно называть \emph{моделью}, а не сырой корреляцией, лишь когда \(M_t\) --- минимально достаточная статистика прошлого для предсказания будущего: критерий модельности --- § 2.3); \(I_{\text{pred}}\) --- \textbf{предсказательность}: какая часть этого знания относится к будущему среды на горизонте \(\tau\). Узкий проксический оператор \(M_t\) --- функциональный канал, а не полная взаимная информация \(I(X; X_E)\) по всем степеням свободы: последняя была бы на много порядков больше за счёт \textbf{нефункциональных сопряжений} со средой (тепловой обмен, механический контакт, диффузия), статистически связывающих \(X\) и \(X_E\), но не несущих predictive information о будущем её поведении. Предсказательная составляющая отличается также от классической predictive information Биалека--Неменмана--Тишби (Bialek et al.~2001), определённой как взаимная информация между прошлым и будущим одного процесса, \(I(X_E^{[t-\tau,\,t]};\, X_E^{[t,\,t+\tau]})\); связь между ними даётся неравенством об обработке данных (см. ниже) и обсуждается далее как мост к рамке Биалека.

Центральная величина статьи --- \textbf{предсказательная эффективность самомоделирования} --- определяется как отношение этих двух величин:

\begin{equation}\eta_v(t,\tau) := \frac{I_{\text{pred}}(t,\tau)}{I_{\text{mem}}(t)} = 1 - \nu(t,\tau), \qquad
  \nu(t,\tau) := \frac{I_{\text{mem}}(t) - I_{\text{pred}}(t,\tau)}{I_{\text{mem}}(t)}, \tag{1}\end{equation}

где \(\nu\) --- \textbf{информационная ностальгия}, доля удерживаемой памяти, переставшая быть предсказательной. \emph{То есть:} \(\eta_v\) --- не «биты предсказания на джоуль», а доля того, что система помнит, которая реально относится к будущему; \(\nu\) --- дополнительная к ней доля чистого балласта.

Граница этой величины --- теорема, а не соглашение о нормировке. Память формируется из прошлого с внутренним шумом, \(M_t = f(X_E^{\le t}, \zeta_v)\), причём шум не несёт информации о будущем среды сверх того, что уже содержится в прошлом: \(\zeta_v \perp X_E^{[t,\,t+\tau]} \mid X_E^{\le t}\) (именно эта условная независимость, а не более сильное требование «отсутствия обратной связи», --- точная нужная посылка; разведение этих двух условий --- supplementary § S2.1). Тогда цепь \(M_t \to X_E^{\le t} \to X_E^{[t,\,t+\tau]}\) марковская, и по неравенству об обработке данных (data processing inequality, DPI; информация о случайной величине не возрастает при детерминированной обработке) \(I_{\text{pred}} \le I_{\text{mem}}\). Отсюда

\[\eta_v \in [0, 1], \qquad \nu \in [0, 1] \qquad \text{— по построению (DPI).}\]

\emph{То есть:} нормировка \([0, 1]\) --- следствие информационного неравенства.

Граница \(\eta_v \le 1\) выведена для \textbf{пассивного} самомоделирования, в котором марковость цепи \(M_t \to X_E^{\le t} \to X_E^{[t,\,t+\tau]}\) держится. Активный случай (свободно плывущая \emph{E. coli}, агенты active inference § 7), в котором \(M_t\) каузально гонит будущее среды через действие и тем самым рвёт цепь DPI, остаётся \textbf{открытым пунктом}: conditioning на фиксированную политику марковости не восстанавливает (действие \(a = \pi(M_t)\) делает \(M_t\) зависимым от будущего среды), и граница требует перехода к \textbf{направленной} (transfer) информации \(T_{M \to X_E}\) --- сколько собственное состояние \(M\) каузально привносит в будущее среды сверх её прошлого --- feedback-обобщение через бонд Sagawa--Ueda (2010); постановка --- § 2.7 ниже.

\emph{Интервенционная (do) легитимность пассивного измерения.} Открытость активного случая не лишает смысла операциональное измерение \(\eta_v\): иммобилизованный протокол § 3 фиксирует не «упрощённую», а строго определённую величину, поскольку реализует \textbf{do-интервенцию} в смысле Пирла (Pearl 2009). Экзогенный зажим действия \(\mathrm{do}(a)\) --- контролируемое возмущение, фиксирующее \(a\) внешней рукой, --- разрывает причинную стрелку \(M_t \to a \to X_E\) и восстанавливает марковость \(M_t \to X_E^{\le t} \to X_E^{[t,\,t+\tau]}\); по DPI

\begin{equation}\eta_v^{\mathrm{do}} := \frac{I\!\left(M_t;\; X_E^{[t,\,t+\tau]} \,\middle|\, \mathrm{do}(a)\right)}{I_{\text{mem}}}
  \in [0, 1]. \tag{1b}\end{equation}

\emph{То есть:} под интервенцией, зажимающей действие экзогенно (а \textbf{не} conditioning на наблюдаемую политику --- это разные операции: пассивное обусловливание оставляет путь \(X_E^{\text{прош}} \to M_t \to a \to X_E^{\text{буд}}\) открытым, тогда как \(\mathrm{do}(a)\) обрезает входящую в \(a\) стрелку из \(M_t\)), петля снова пассивна и отношение корректно лежит в \([0, 1]\). Знаменатель при этом остаётся наблюдательным \(I_{\text{mem}}\), а не интервенционным: зажим \(\mathrm{do}(a)\) в настоящем и будущем не меняет \textbf{уже сформированную} корреляцию \(M_t\) с прошлым среды (прошлый сегмент \(X_E^{\le t}\), на котором оценивается \(I_{\text{mem}}\), формируется \textbf{до} зажима), поэтому \(I\!\left(M_t;\, X_E^{\le t}\,\middle|\,\mathrm{do}(a)\right) = I_{\text{mem}}\), и DPI под интервенцией замыкает \(\eta_v^{\mathrm{do}} \le 1\) с числителем и знаменателем в \textbf{одном} интервенционном режиме (иммобилизованный протокол § 3 фиксирует и тот, и другой). Для \emph{E. coli} физической реализацией \(\mathrm{do}(a)\) служит \textbf{иммобилизация}: пришпиленная к подложке клетка не плывёт, её состояние не движет концентрацию лиганда, и стрелка \(M \to a \to X_E\) разорвана. Поэтому иммобилизованный FRET-протокол Sourjik--Berg (§ 3) --- \textbf{не упрощение, а операционная реализация} \(\mathrm{do}(a)\), делающая \(\eta_v\) корректной \([0, 1]\)-величиной; свободно плывущая клетка измеряет иную, активную величину, не ограниченную единицей. При этом \(\eta_v^{\mathrm{do}}\) --- \textbf{не новая теорема}, а стандартная do-редукция активного контура к пассивному, фиксирующая лишь точный смысл легитимности пассивного измерения. Не следует смешивать \(\eta_v^{\mathrm{do}}\) с термодинамическими КПД той же feedback-литературы: ограниченный единицей feedback-КПД существует --- термодинамический \(\eta_X = \dot I/\dot S_r^X \le 1\) (Horowitz and Esposito 2014, ур. (18)) и сенсорная ёмкость \(C = \dot I/T_{y\to x} \in [0, 1]\) (Hartich, Barato, and Seifert 2016), определения --- § 2.7, --- но это \textbf{термодинамическая/сенсорная} величина (поток информации против производства энтропии, точность сенсора), \textbf{не} предсказательная свежесть \(\eta_v\); отождествлять их нельзя.

\emph{Согласование \(S\) и \(\mathrm{do}(a)\).} Зажим \(\mathrm{do}(a)\), приостанавливающий \textbf{action-ветвь} петли возврата ошибки (ii) на время измерения \(\eta_v^{\mathrm{do}}\), не вступает в конфликт с присвоением \(S = 1\). Предикат самооплаты \(S\) требует, чтобы удерживаемая моделью информация и оплачивающая её удержание диссипация принадлежали одной физической границе системы --- само-питание (условие (i)) и возврат ошибки модели в собственный контур (условие (ii)); полное определение --- § 2.7. Это \textbf{диспозициональный} предикат \emph{нативного} (свободно плывущего) организма, устанавливаемый \textbf{отдельной} интервенцией --- counterfactual-ablation (§ 2.2), --- а не do-зажимом. У петли (ii) \textbf{две ветви} --- action-канал (сенсомоторный возврат ошибки) и диссипативная Still-ветвь; \(\mathrm{do}(a)\) приостанавливает \textbf{лишь action-ветвь} на время оценки \(\eta_v^{\mathrm{do}}\) и не отменяет диспозициональный \(S\); то есть \(S\) и \(\eta_v^{\mathrm{do}}\) оценивают одну систему \textbf{разными} протоколами, и кажущегося конфликта «зажав \(a\), как утверждается \(S = 1\)?» нет --- вердикт \(V(X) := S \wedge (\eta_v > 0)\) соединяет две разные функции системы, измеряемые двумя разными интервенциями. Для иммобилизованного препарата \textbf{диссипативная ветвь (ii) остаётся активной} и держит \(S = 1\): инстанцируясь не через перерезанный action-канал, а через \textbf{диссипативный механизм} Still --- искажение модели → рост непредсказательной корреляции \(I_{\text{nonpred}}\) → больший Still-налог, оплачиваемый \textbf{собственным АТФ} клетки (условие (ii) § 2.7 уже допускает «рост обязательной диссипации \textbf{или} распад \(X\)»).

Перед привязкой \(\eta_v\) к фиксированному горизонту вводится характерное окно системы. Горизонт \(\tau_d := I_{\text{pred}}/\dot I_{\text{pred}}\) --- memory turnover time, отношение запаса predictive information к темпу её поступления; он определён для \textbf{немарковских} (память-несущих) сред. Для марковских сред \(\dot I_{\text{pred}} = 0\) (горизонт не добавляет предсказательной информации сверх одного шага), и характерный горизонт задаётся временем корреляции среды \(\tau_E\). Эмпирические окна --- время адаптации сигнального контура (секунды--минуты, внутри жизни одной клетки) для бактерии, год для города, время отклика биогеохимического цикла (год--десятилетие) для биосферы --- суть методологические прокси для \(\tau_d\), чья согласованность с теоретическим горизонтом есть отдельный калибровочный вопрос (§ 3.5). Ниже всюду \(\eta_v := \eta_v(t, \tau_d)\).

Величина \(\eta_v\) имеет статус \textbf{внутриклассового грейда} свежести self-model --- она измеряет, насколько устарела модель данной системы, --- \textbf{а не} межсистемной абсолютной шкалы. Отношение \(I_{\text{pred}}/I_{\text{mem}}\) штрафует за объём памяти, поэтому им нельзя напрямую сравнивать «бактерию против мозга»: межсистемная демаркация (синхронной) витальности есть работа предиката \(S\) (§ 2.7): бинарный тест «витален или нет» делает предикат, а численное значение \(\eta_v\) внутри положительной области даёт сравнительную меру эффективности удержания модели. Почему именно самооплата (а не autopoiesis, RAF или организационное замыкание) служит критерием принадлежности и почему \(S\) --- осознанная онтологическая посылка, а не произвольное определение, разбирается в § 2.7 (эпистемический статус): информационная теория отвечает лишь за \textbf{эффективность} модели и сама не определяет, \textbf{кто оплачивает} её удержание; предикат \(S\) добавляет именно этот физический критерий принадлежности, с позитивным аргументом необходимости (внутрисистемная отбраковка ошибок) и комплементарностью --- не конкуренцией --- с линией организационного замыкания.

\emph{Об измеримости знаменателя.} Знаменатель \(I_{\text{mem}} = I(M_t; X_E^{\le t})\) измерим из временных рядов так же, как фактическая диссипация \(E_{\text{actual}}\) --- величина, доступная через калориметрию: для хемотаксической петли \emph{E. coli} в пассивном (иммобилизованном) FRET-протоколе Sourjik--Berg он зависит лишь от ковариаций метилирования (латентная переменная, восстанавливаемая из FRET-считывания активности киназы) и лиганда (микрофлюидика); в гауссовом приближении линейной модели оценивается через logdet этих ковариаций (полный proof-of-measurability и worked-example --- § 3).

\textbf{Термодинамический якорь (Still 2012).} Балласт памяти не бесплатен: чтобы система поддерживала корреляцию с уже-неактуальным прошлым, термостат обязан получить тепло. Still доказал точную нижнюю цену этого балласта, и она служит в настоящей рамке \textbf{отдельным термодинамическим якорем}, не входящим в определение (1). На шаге обновления памяти \(t \to t+1\) для мгновенных величин \(I_{\text{mem}}^{\text{инст}} = I(M_t; X_{E,t})\) и \(I_{\text{pred}}^{\text{инст}} = I(M_t; X_{E,t+1})\) Still доказывает \textbf{равенство} (его уравнение (14)), \(\langle W_{\text{diss}}^{t\to t+1}\rangle = k_B T\,(I_{\text{mem}}^{\text{инст}} - I_{\text{pred}}^{\text{инст}})\); для полного протокола (работа и последующая релаксация) неотрицательный релаксационный член превращает его в \textbf{нижнюю границу} (bound) Still --- далее бонд (2) (его уравнение (18)). Still-цена считается по-шагово --- каждый шаг обновления памяти оплачивается отдельно, и за окно эти платы складываются; поэтому различаем мгновенную ностальгию \(\nu^{\text{инст}}\) (плата одного шага) и горизонтную \(\nu(t,\tau)\) (доля балласта за всё окно):

\begin{equation}\langle W_{\text{diss}} \rangle \;\ge\; k_B T\,\big(I_{\text{mem}}^{\text{инст}} - I_{\text{pred}}^{\text{инст}}\big)
  \;=\; k_B T\,\nu^{\text{инст}}\,I_{\text{mem}}^{\text{инст}}. \tag{2}\end{equation}

\emph{То есть:} каждый бит мгновенной ностальгии стоит системе не менее \(k_B T \ln 2\) рассеянного тепла --- это строгий термодинамический смысл удержания устаревающей модели. Две оговорки фиксируют точный статус бонда. Во-первых, (2) несёт \textbf{мгновенную} ностальгию \(\nu^{\text{инст}} = 1 - I_{\text{pred}}^{\text{инст}}/I_{\text{mem}}^{\text{инст}}\) (слой Still), не горизонтную \(\nu(t,\tau)\) из (1); кумулятивная диссипация за окно есть \textbf{сумма} пошаговых вкладов, а не произведение \(k_B T\,\nu^{\text{кум}} I_{\text{mem}}^{\text{кум}}\) (пошаговая структура диссипации и её связь с горизонтной \(\nu(t,\tau)\) --- supplementary § S2.1.1). Во-вторых, \(k_B T\) входит сюда как \textbf{коэффициент границы из теоремы} --- множитель в неравенстве, --- а \textbf{не} как курс обмена энергии на биты; единственное законное вхождение ландауэровского масштаба \(k_B T \ln 2\) --- цена стирания одного бита внутри (2). Этот якорь проясняет связь витальности с необратимостью --- но в тонком смысле, который легко исказить инверсией. Живой аппарат селектируется \textbf{к} ландауэровскому пределу \textbf{пошагово} (Still 2012): биомолекулярные машины почти-обратимы за шаг (порядка \(20\,k_B T\) на операцию --- Bennett 1982), и эта пошаговая эффективность есть \textbf{признак} витальности, не её отрицание. Исключается лишь \textbf{глобальная} обратимость всей операции: несущее условие (ii) предиката \(S\) --- это коррекция ошибки модели, а коррекция (кинетическая корректура, proofreading --- Hopfield 1974; Ninio 1975) термодинамически необходимо диссипативна (Bennett 1982; Landauer 1961: вечно работающая полезная машина не может хранить всю историю), так что строго обратимый самомоделирующий контур \textbf{не способен} корректировать ошибку и потому не удовлетворяет (ii) --- обе её ветви, «рост обязательной диссипации» и «распад \(X\)» (сам необратимый), восходят к одному запрету. То есть \(S = 0\) здесь следует \textbf{по построению предиката}, а не из обнулённой магнитуды диссипации (информационная \(\eta_v\) при этом выживает --- граница из DPI, не из стирания); в этом ограниченном смысле витальность правдоподобно конститутивно необратима --- тезис, формальное проведение которого остаётся открытой задачей (строгий маршрут и различение пошагово/глобально --- supplementary § S2.1.1; область определения \(T_v\)-якоря --- § S5.1).

Связь predictive information с диссипацией идёт через \textbf{изменение энтропии и взаимной информации} (бонд Still (2); обобщённый принцип Ландауэра, Sagawa 2014, P03025), а не через пересчёт свободной энергии в биты. В определении (1) масштаба \(k_B T \ln 2\) нет; линия «Landauer» появляется законно --- через бонд (2), где и только где входит \(k_B T\). При \(\dot I_{\text{mem}} \to 0\) (память стабилизирована, \(I_{\text{mem}} \approx \text{const}\)) и стационарном дрейфе ностальгия \(\nu\) выходит на постоянный уровень, а кумулятивная диссипация растёт линейно с положительной удельной скоростью \(\ge k_B T \varepsilon\) --- строгий термодинамический смысл неустранимости сноса в стационаре (supplementary § S2.1.1); нестационарный режим разбирается в сопровождающей работе (Andriishin 2026).

\subsection{Граница системы как методологическая конвенция}\label{ux433ux440ux430ux43dux438ux446ux430-ux441ux438ux441ux442ux435ux43cux44b-ux43aux430ux43a-ux43cux435ux442ux43eux434ux43eux43bux43eux433ux438ux447ux435ux441ux43aux430ux44f-ux43aux43eux43dux432ux435ux43dux446ux438ux44f}

Шкала \(\eta_v\) зависит от выбора границы измеряемой системы. Это структурное свойство, не дефект. Для крупной языковой модели при строгой границе «веса в момент инференса» (прямой проход без сохранения состояния) одновременно проваливаются два условия: собственная диссипация \(E_{\text{actual}}^{\text{own}} \to 0\) (электросеть внешняя) и канал возврата ошибки в работающий контур отсутствует. Шкала \(\eta_v\) при такой границе \textbf{не определена} через одновременный провал обоих условий, не «равна нулю через числитель» (формальный счёт через ёмкости \(T_v\), \(C_v\), \(S_v\) --- § 2.3; различение строгого/операционального прочтения \(I_v^{\text{int}}\) против \(I_v^{\text{op}}\) --- § 2.5). При границе, включающей дата-центр и корпорацию, \(\eta_v > 0\), но измеряемый объект --- уже не модель.

Выбор границы методологически дисциплинирован. Граница проводится там, где замкнута петля принятия решений, оплачиваемых собственным распадом системы. Концептуальный предшественник --- \emph{epistemic cut} и \emph{semantic closure} Паттее (Pattee 1982, 2001): граница между физико-динамическим уровнем и символическим уровнем, в которой система замыкает в себе обработку собственных символов. Бактериальная клетка --- естественная граница, поскольку искажение её внутренней модели градиента ведёт к распаду именно этой клетки. Город --- поскольку искажение его модели потоков возвращается кризисом именно к этой административной единице. Тождественность вопроса определяет границу; выбор границы задаёт измеряемое.

Конвенциональность границы продуктивна. Возражение «при правильном выборе границы любая система окажется витальной» парируется так: выбор границы есть выбор того, чью витальность измеряют, и расширение границы расширяет ответственный субъект, а не растворяет диагноз.

\emph{Граница и Markov blanket.} Markov blanket системы \(X\) (Friston 2019) --- статистически разделяющий слой сенсорных и активных состояний. Собственная конструкция self-payment-границы (§ 2.2) \textbf{не наследует} трудностей реалистской интерпретации Markov blanket, диагностированных в BDDB-2022 (Bruineberg et al.~2022) и Aguilera et al.~(2022). Граница в нашей рамке задана \textbf{термодинамической принадлежностью бюджета} \(E_{\text{actual}}\), а не статистическим разделением распределения по Pearl- или Friston-blanket; поэтому ортогональна Pearl/Friston-blanket дилемме. Смешение двух понятий приводит к категориальной ошибке, где информационное разделение принимается за энергетическую самооплату.

\emph{Counterfactual ablation.} Дисциплина выбора границы операционализируется интервенционной процедурой в духе Pearl-овой причинности (терминологически --- counterfactual ablation в смысле mechanistic interpretability: удалить компонент, измерить эффект; не путать с Reichenbach screening-off). Для кандидата \(c\) на включение в границу \(X\): \(c\) вычитается из физической реализации; если изменение скорости истощения свободной энергии \(-dF_X/dt\) за время отклика \(< \tau_d\) значимо, \(c\) принадлежит границе. Количественный порог: \(c\) принадлежит границе при \(|\Delta(-dF_X/dt)| / |-dF_X/dt|_{\text{base}} \ge \delta\), где \(\delta \in [0{,}05;\, 0{,}1]\) --- методологическая конвенция, аналогичная порогу значимости в регрессионном анализе. Для систем с плавно распределённой принадлежностью (паразит-симбионт-хозяин, корпорация-сотрудники) процедура возвращает шкалу степени принадлежности --- степень принадлежности \(w_c = |\Delta(-dF_X/dt)| / \sum_i |\Delta(-dF_{X,i}/dt)| \in [0, 1]\).

Существенная оговорка о \emph{том, что} измеряет интервенция: эффект вычитания \(c\) берётся как вклад в \textbf{замкнутую корректирующую петлю регулирования} --- топологически, по тому, лежит ли \(c\) на пути «возмущение регулируемой переменной → корректирующее действие» с откликом за время \(< \tau_d\), --- а \textbf{не} как сырой энергопоток сам по себе. Без этого уточнения чистый источник питания с наибольшим \(|\Delta(-dF_X/dt)|\) (электросеть термостата, чьё удаление обнуляет всю диссипацию) попал бы \textbf{внутрь} границы; топологическое различение петли оставляет его снаружи --- он \emph{питает} диссипацию, но не входит в контур, отвечающий на возмущение регулируемой переменной (биметалл и реле входят, электросеть и нагреватель --- нет).

Процедура опирается на pre-theoretic candidate boundary --- минимальную начальную интуицию (физическая связность, биологическая клетка как единица, инженерная сборка как единица). Counterfactual-вычитание уточняет границу через измерение \textbf{относительного} вклада каждого компонента в полную диссипацию связной системы; последовательное уточнение, устойчивое к перестановкам вычитаемых элементов, даёт обоснованную границу. Это не полная элиминация конвенциональности (для систем без естественной pre-theoretic кандидатуры --- корпорация, биосфера --- процедура остаётся открытой), но операциональная дисциплина выбора в большинстве физических и биологических случаев.

\emph{Членство-в-границе \ensuremath{\ne} самооплата.} Две операции разведены и применяются строго по порядку. Принадлежность кандидата \(c\) границе решается \textbf{только} вкладом \(c\) в петлю регулирования (\(|\Delta(-dF_X/dt)| \ge \delta\) под возмущением регулируемой переменной --- разрывом петли при вычитании \(c\), топологически, а не по сырому энергопотоку) --- до и независимо от вопроса о плательщике (i); предикат самооплаты \(S\) (§ 2.7) оценивается \textbf{затем}, уже на зафиксированной границе. Порядок «граница → потом \(S\)» разрывает именно порочную петлю «границу подбирают, чтобы \(S\) вышел нужным»: для термостата ablation относит к границе биметалл и реле (их вычитание рвёт контур регулирования), а нагреватель и электросеть оставляет снаружи --- и это решается без всякой апелляции к самооплате; лишь \textbf{после} фиксации границы выясняется \(S = 0\), ибо обязательную диссипацию несёт сеть, не биметалл (полный разбор --- § 4.4). Остаточная конвенциональность (выбор кандидатной границы, порог \(\delta\)) при этом сохраняется (§ 2.7): механизм дисциплинируется, но не устраняется.

Применение к LLM --- § 3.4 и Supplementary § S3.4; к термостату --- § 4.4 и Supplementary § S4.4.

\subsection{Триада инвариантов}\label{ux442ux440ux438ux430ux434ux430-ux438ux43dux432ux430ux440ux438ux430ux43dux442ux43eux432}

Положительность \(\eta_v\) требует одновременного присутствия трёх ёмкостей.

\emph{Термодинамическая ёмкость \(T_v\)} --- цена удержания модели против эрозии. Необратимое стирание бита физически стоит не менее \(k_B T \ln 2\) свободной энергии (Landauer 1961); поток эксергии --- плата против шума. Стандартный прокси --- удельный поток эксергии \(P_D \cdot \eta_{\text{ex}}\), нормированный на массу или объём; в традиции Чейссона (Chaisson 2001) измерим от молекул до галактик.

\emph{Информационная ёмкость \(C_v\)} --- качество компрессии: адекватность модели среде и обновляемость. Прокси разделяются на два класса. Первый --- прокси скорости энтропии \(h_\mu\) (нормированная Lempel--Ziv-сложность, перестановочная энтропия Бандта--Помпе) --- операционализирует темп новизны среды, не предсказательную информацию. Второй --- прокси \(I_{\text{pred}}\):

\begin{itemize}
\tightlist
\item
  excess entropy \(\mathbf{E} := \lim_{L\to\infty} I(X^{[t-L,\,t]};\, X^{[t,\,t+L]})\) --- асимптотическое значение predictive information в формулировке Биалека--Неменмана--Тишби в пределе бесконечных окон; на конечных окнах операционализируется через block-entropy оценки \(H(L)\) (Crutchfield and Feldman 2003);
\item
  statistical complexity \(C_\mu\) (вычислительная механика Crutchfield; Shalizi and Crutchfield 2001) --- энтропия минимальной достаточной статистики для предсказания, удовлетворяет \(C_\mu \ge \mathbf{E}\) и служит дополнительной мерой структурной памяти;
\item
  PCI (perturbational complexity index) Казали (Casali et al.~2013) --- оценка нелинейной интеграции на корково-биологических данных. Параллельная программа Адами (Adami 2002) определяет physical complexity как взаимную информацию между геномом и окружением --- биологически-конкретный частный случай «качества модели среды», концептуально родственный \(C_v\). Различение классов принципиально для последовательной операционализации \(\eta_v\). В стационарном режиме настоящей работы оба класса служат прокси одного \(C_v\). Тождество \(\eta_v = C_v^{\text{predictive}}/C_v^{\text{static}}\) корректно лишь внутри \textbf{второго} класса (вычислительная механика), где обе величины --- биты удержанной структуры: \(C_v^{\text{static}}\) отождествляется со statistical complexity \(C_\mu\) (структурная память, прокси \(I_{\text{mem}}\)), а \(C_v^{\text{predictive}}\) --- с excess entropy \(\mathbf{E}\) (предсказательная часть, прокси \(I_{\text{pred}}\)). Теорема \(C_\mu \ge \mathbf{E}\) тогда даёт \(\eta_v = \mathbf{E}/C_\mu \in [0, 1]\) согласованно с DPI. Прокси \textbf{первого} класса (\(h_\mu\), Lempel--Ziv) измеряют темп новизны среды --- это \textbf{отдельный} индекс, не знаменатель \(\eta_v\). Тем самым центральная величина § 2.1 получает триадную форму: предсказательная эффективность есть доля удерживаемой структурной памяти (\(C_\mu\)), остающаяся предсказательной (\(\mathbf{E}\)) --- \(\eta_v = I_{\text{pred}}/I_{\text{mem}} = 1 - \nu = C_v^{\text{predictive}}/C_v^{\text{static}}\); тождество точно, когда \(M_t\) --- каузально-достаточное (минимально достаточное) представление, и есть прокси-приближение для произвольного функционального канала \(M_t\). Именно эта оговорка несёт \textbf{критерий модельности}, отделяющий \emph{модель} от сырой памяти: удержанное есть модель среды, а не произвольная корреляция, ровно тогда, когда \(M_t\) --- каузально-достаточная (минимально достаточная) статистика прошлого для предсказания будущего, то есть когда \(I_{\text{mem}}\) отождествимо с \(C_\mu\) (энтропией каузальных состояний, Shalizi and Crutchfield 2001), а не с раздутой нефункциональными сопряжениями сырой \(I(X; X_E)\) (§ 2.1). Тогда \(\eta_v = \mathbf{E}/C_\mu\) измеряет ровно ту долю минимально-достаточной памяти, что несёт предсказательную структуру (\(\mathbf{E}\)); это и есть основание называть \(I_{\text{mem}}\) \textbf{моделью}, а не просто памятью --- и правомерность тем полнее, чем ближе функциональный канал \(M_t\) к минимальной достаточной статистике. В нестационарном расширении (Andriishin, сопровождающая работа) расхождение \(C_v^{\text{static}}\) и \(C_v^{\text{predictive}}\) служит диагностическим сигналом информационной ностальгии (полное соотношение \(C_v^{\text{predictive}} = (1-\nu)\,C_v^{\text{static}}\) --- § 6).
\end{itemize}

\emph{Топологическая ёмкость \(S_v\)} --- архитектура связности петли возврата ошибки. Стандартные прокси --- модулярность Ньюмана-\(Q\) (Newman and Girvan 2004) с ориентировочным порогом \(Q > 0{,}3\) как условным ориентиром, требующим калибровки нуль-моделью (см. ниже про resolution limit, а не жёсткий критерий), безмасштабность \(\gamma \in [2, 3]\) для распределения степеней (Barabási and Albert 1999; Newman 2010 обзор; критика универсальности в Broido and Clauset 2019), иерархический коэффициент кластеризации \(C(k) \sim k^{-1}\) (Ravasz and Barabási 2003). Алгебраическая связность \(\lambda_2\) Fiedler (1973), индексирующая глобальную диффузионную связность, --- \emph{вторичный} прокси с обратной интерпретацией: для модулярных архитектур \(\lambda_2\) низкая (Mohar 1991), поэтому она используется как индикатор узких мест структуры между модулями, не как со-направленная с \(Q\) мера. Дополнительный класс прокси --- структурная управляемость (Liu, Slotine, and Barabási 2011) --- независимый от модулярности срез топологии петли возврата ошибки.

Эффективность иерархически-модулярной архитектуры для удержания predictive information над структурированной средой формализуется через минимум описательной длины (Mengistu et al.~2016): иерархия топологии сжимает описание адаптивной политики на величину, пропорциональную глубине иерархии среды. Иерархическая структура валидируется через предсказательную силу разложения на missing-links (Clauset, Moore, and Newman 2008), что эмпирически связывает \(S_v\)-прокси с \(C_v\)-критериями и предотвращает их вырождение в одну меру. Эмпирический коррелят --- нейросетевой connectome (Bassett and Sporns 2017; Lynn and Bassett 2019), демонстрирующий иерархически-модулярную организацию как селекционный коррелят вычислительной эффективности. Два крайних случая обнуляют \(S_v\) конструктивно (§ 5, Блок 1 «Прямые опровержения», п. 3). Случайный граф Erdős--Rényi --- отсутствие модулярной декомпозиции; избыточные короткие связи не дают сжатия адаптивной политики, несмотря на small-world diameter (Watts and Strogatz 1998). Регулярная решётка --- отсутствие межмасштабной связи, проигрыш на крупномасштабное распространение коррекции ошибки.

Эмпирический статус scale-free как универсального паттерна оспорен: на корпусе \textasciitilde1000 эмпирических сетей строгое power-law с \(\gamma \in (2, 3)\), проверяемое по статистической методике Clauset et al.~(2009), подтверждается в меньшинстве случаев (Broido and Clauset 2019; Holme 2019). Степенное распределение --- один из режимов требуемой архитектуры, не единственный; релевантно более общее условие отсутствия характерного масштаба и наличия модулярно-иерархической организации. \(Q\)-метрика Newman зависит от алгоритма (greedy Newman (2004), Louvain (Blondel et al.~2008), Leiden (Traag et al.~2019)); resolution limit problem (Fortunato and Barthélemy 2007) требует калибровки порога через нуль-модель. В настоящей работе \(Q\) считается алгоритмом Leiden; численный расчёт для хемотактической сети \emph{E. coli} --- supplementary § S2.4 (\(Q \sim 0{,}5\)). Эволюционное происхождение иерархической модулярности (введённой выше через \(C(k) \sim k^{-1}\) Ravasz--Barabási) --- давление отбора на минимизацию длины связей (Mengistu et al.~2016) и изменчивость среды (Kashtan and Alon 2005).

Триада не сложена из независимых дисциплин --- она вынуждена природой удержания сжатия во времени. Обнуление любой ёмкости обнуляет витальность по любой шкале: без оплаты нет ставки --- собственной цены системы за качество удерживаемой модели, делающей её ошибки значимыми для самой системы (skin-in-the-game, § 2.7), --- без качества нет модели, без архитектуры нет петли возврата.

\subsection{Нормировка ёмкостей по референтной когорте}\label{ux43dux43eux440ux43cux438ux440ux43eux432ux43aux430-ux451ux43cux43aux43eux441ux442ux435ux439-ux43fux43e-ux440ux435ux444ux435ux440ux435ux43dux442ux43dux43eux439-ux43aux43eux433ux43eux440ux442ux435}

Прокси в прямой форме не лежат на \([0, 1]\) и не подставимы в композит \(I_v\) без нормировки. Удельный поток эксергии для \(T_v\) имеет размерность Вт·кг\(^{-1}\) и неограничен сверху. Excess entropy и \(C_\mu\) для \(C_v\) --- биты, неограниченные сверху. Модулярность Newman для \(S_v\) лежит в \([-1/2, 1]\); иерархичность Ravasz--Barabási --- степенной показатель. \(\lambda_2\) для связного графа \(\in (0, n]\) (для несвязного \(\lambda_2 = 0\)); как вторичный прокси не используется в составе \(S_v\)-нормировки. Без явной нормировки формула (2b) и оговорка «\(T_v, C_v, S_v \in [0, 1]\)» формально некорректны.

В настоящей работе принята перцентильная нормализация по референтной когорте. Для каждой ёмкости фиксируется опорная когорта систем того же структурного класса; нормированная величина --- перцентиль в распределении когорты:

\begin{equation}\tilde{T}_v(X) = F_{T_v}^{\text{cohort}}\bigl(T_v(X)\bigr) \in [0, 1], \tag{2a}\end{equation}

где \(F_{T_v}^{\text{cohort}}\) --- эмпирическая функция распределения прокси по когорте; аналогично для \(\tilde{C}_v\), \(\tilde{S}_v\). Нулевой перцентиль --- слабейший представитель, единичный --- сильнейший. Стандартный приём социальных индикаторов: Human Development Index (United Nations Development Programme 2022) нормирует свои компоненты так же. Конкретные когорты для шести камертонов и обработка краевых случаев (биосфера без когорты; кристалл и ураган с конструктивным нулём) --- Supplementary § S2.4.

Перцентильная процедура несёт три ограничения. Первое: выбор когорты --- конвенция. При этом устойчивость различия по знаку между двумя системами одной когорты не зависит от выбора когорты при достаточной её ширине; именно эта устойчивость по знаку --- содержательное утверждение § 3, а не конкретные числа. Второе: процедура теряет информацию об абсолютных значениях прокси и работает только для сравнения внутри класса. Между классами численные \(I_v\) напрямую несопоставимы. Третье: процедура неприменима к системам без естественной когорты, где требуется нормировка по теоретическому верхнему пределу. Сигмоидная альтернатива \(\tilde{T}_v = T_v/(T_v + T_v^{\text{ref}})\) --- Supplementary § S2.4.

Далее в тексте \(T_v, C_v, S_v\) обозначают нормированные величины \(\tilde T_v, \tilde C_v, \tilde S_v\), если явно не указано иное; тильда опускается ради читаемости.

\subsection{\texorpdfstring{Композит \(I_v\) как индекс структурной готовности}{Композит I\_v как индекс структурной готовности}}\label{ux43aux43eux43cux43fux43eux437ux438ux442-i_v-ux43aux430ux43a-ux438ux43dux434ux435ux43aux441-ux441ux442ux440ux443ux43aux442ux443ux440ux43dux43eux439-ux433ux43eux442ux43eux432ux43dux43eux441ux442ux438}

Геометрическое среднее трёх ёмкостей

\begin{equation}I_v = \sqrt[3]{\tilde{T}_v \cdot \tilde{C}_v \cdot \tilde{S}_v} \in [0, 1] \tag{2b}\end{equation}

определяется как структурный композит. Геометрическое усреднение выбрано из трёх соображений: остаётся в \([0, 1]\); симметрично по перестановкам; обнуляется при обнулении любой ёмкости (выпадение одной ёмкости обрушивает витальность целиком). Аддитивное усреднение этого свойства не сохраняет --- сумма маскировала бы провал. Веса \(w_T = w_C = w_S = 1\) --- конвенция; для неравных весов (при \(w_i > 0\)) \(I_v^{(w)} = (\tilde{T}_v^{w_T} \cdot \tilde{C}_v^{w_C} \cdot \tilde{S}_v^{w_S})^{1/(w_T + w_C + w_S)}\). Анализ чувствительности к весам не проводится, так как \(I_v\) есть методологический индекс структурной готовности к витальности, не самостоятельная витальная шкала: бинарный тест «жив или нет» делает \(\eta_v\), не \(I_v\). Положительный \(I_v\) --- необходимое условие положительного \(\eta_v\); достаточным он не является. Через стандартные прокси (Lempel--Ziv, эксергия, сетевые метрики) \(I_v\) даёт ненулевые значения у любой структурно сложной диссипативной системы --- Солнца, урагана, чёрной дыры, LLM. По позиционному месту \(I_v\) совпадает с assembly index, морфодинамикой Дикона, \(\Phi\) Тонони: композитный индекс сложности, необходимый, но не достаточный.

Возражение к статусу \(I_v\) могло бы состоять в том, что композит трёх необходимых условий естественнее было бы интерпретировать как самостоятельную шкалу. Однако \(I_v\) через любой разумный набор индустриальных прокси даёт устойчиво положительные значения у заведомо невитальных систем. Например, \(I_v^{\text{op}}\) Солнца положителен через любой стандартный прокси. Нормированная Lempel--Ziv-сложность X-flux GOES за 11-летний цикл даёт \(\sim 0{,}5\)-- \(0{,}7\) от верхнего перцентиля когорты звёзд главной последовательности; иерархичность сети активных областей по SDO/HMI (Gheibi et al.~2017) --- \(\sim 0{,}4\). Конкретное число зависит от датасета и когорты, но устойчиво \(\in (0{,}3;\, 0{,}7)\). Поэтому \(I_v\) остаётся в роли \emph{методологического} индекса; формально это закрепляется различием в двух способах прочтения: \(I_v^{\text{op}}\) (operational reading) --- через стандартные прокси на наблюдаемых данных --- ненулевое у любой сложной диссипативной системы. \(I_v^{\text{int}}\) (interpretive / strict reading) --- со строгой проверкой \(T_v\) как собственной диссипации, \(C_v\) как модели среды в техническом смысле, \(S_v\) как архитектуры замкнутой петли --- равно нулю у любой невитальной системы по построению. Расхождение двух прочтений у конкретной системы --- диагностический сигнал: система структурно сложна, но не витальна, и расхождение указывает ёмкость провала.

\subsection{Двухступенчатая процедура}\label{ux434ux432ux443ux445ux441ux442ux443ux43fux435ux43dux447ux430ux442ux430ux44f-ux43fux440ux43eux446ux435ux434ux443ux440ux430}

Витальность теоретически устанавливается одним тестом --- на положительность \(\eta_v\). Двухступенчатая процедура --- методологическая удобность, не логическая необходимость:

\begin{enumerate}
\def\labelenumi{\arabic{enumi}.}
\tightlist
\item
  \emph{Скрининг \(I_v\).} Дешёвый, через стандартные операциональные прокси; отбраковывает структурно непригодные системы (кристалл --- обнуление \(T_v\); термическое равновесие --- обнуление всех трёх ёмкостей). Вычислительный фильтр, не отдельный тест.
\item
  \emph{Верификация \(\eta_v\).} Дорогая, требует идентификации границы и собственной петли самооплаты; применяется к кандидатам, прошедшим скрининг. Единственный логически содержательный тест.
\end{enumerate}

Бинарность знака \(\eta_v\) не делает численное значение декоративным. Возражение «если демаркацию даёт уже знак, число избыточно --- counterfactual-ablation § 2.2 устанавливает знак без формулы» парируется указанием на нередуцируемое к знаку содержание величины. Во-первых, \(\eta_v\) внутри положительной области --- сравнительная метрика эффективности петли самооплаты: порядок величины \(\eta_v\) у двух систем одного класса несёт содержание сверх общего знака. Во-вторых, величина входит в фальсифицируемое количественное предсказание \(r_{\text{adapt}} = A \cdot \eta_v^{\alpha}\) (§ 5), которое знак не порождает: знак не задаёт ни наклон \(\alpha\), ни проверяемую монотонную зависимость темпа адаптации. Знак \(\eta_v\) отделяет лишь класс без модели (\(\eta_v = 0\)); бинарную демаркацию (синхронной) витальности среди модель-несущих систем несёт предикат \(S\) (§ 2.7), а величина \(\eta_v\) внутри \(S = 1\) --- внутриклассовая метрика и рискованное предсказание; последнее не выводимо из знака.

Структура исходов: прошёл оба порога --- витальные (бактерия, лес, мозг, биосфера, город); прошёл скрининг, провалил верификацию --- структурно богатые, но невитальные (Солнце, ураган, чёрная дыра; LLM при строгой границе «модель как агент» --- провал \(S = 0\)); провалил скрининг --- структурно непригодные (кристалл, термическое равновесие).

Опровержение любой ёмкости обрушивает обе шкалы сразу --- вердикт витальности \(V(X) = S \wedge (\eta_v > 0)\) и структурный скрининг \(I_v\): провал \(C_v\) действует через \(\eta_v\) (через \(I_{\text{mem}}\) или \(I_{\text{pred}}\)), провал \(T_v\) или петли возврата ошибки \(S_v\) --- через предикат \(S\), а любая из трёх обрушивает \(I_v\) через геометрическое среднее. Опровержения двух шкал --- одна и та же диагностика на двух уровнях разрешения, не независимые свидетельства.

Двухступенчатая архитектура несёт методологический риск. Класс «структурно богатые, но невитальные» --- Солнце, ураган, кристалл, LLM, чёрная дыра, реакция Белоусова--Жаботинского --- рискует превратиться в защитный пояс по Лакатошу: любое опровергающее наблюдение поглощается включением в класс. Программа фиксирует минимальную дисциплину: каждое обнуление \(\eta_v\) сопровождается указанием конкретной провалившейся ёмкости. Если класс расширяется без структурного объяснения каждого нового включения, программа сдвигается из прогрессивного в регрессивный режим.

Возможный контр-аргумент: асимметрия двух уровней --- артефакт мягких прокси \(I_v\) (Lempel--Ziv, перестановочная энтропия). Будь \(C_v\) операционализирован строже --- например, через \(\Phi\) Тонони с фильтром на предсказательную силу относительно собственной среды, --- у Солнца получился бы \(C_v \to 0\). Ответ: дело не в выборе одного прокси, а в классе доступных индустриальных метрик. Все прокси сетевой науки, теории информации и неравновесной термодинамики, развитые за последние сорок лет, при применении к сложной диссипативной системе дают ненулевые значения. Прокси, заведомо обнуляющий \(C_v\) Солнца, требует post-hoc калибровки на саму гипотезу витальности --- круг.

\subsection{Самооплата и эпистемический статус}\label{ux441ux430ux43cux43eux43eux43fux43bux430ux442ux430-ux438-ux44dux43fux438ux441ux442ux435ux43cux438ux447ux435ux441ux43aux438ux439-ux441ux442ux430ux442ux443ux441}

Величина \(\eta_v\) говорит, насколько свежа модель, но не отвечает на главный вопрос программы --- \textbf{кто платит} за её удержание. Биметалл в термостате тоже «моделирует» температуру, но диссипацию поддержания оплачивает электросеть; \emph{E. coli} добывает АТФ сама. Витальность наступает тогда, когда система \textbf{сама оплачивает} обязательную --- по Still --- цену удержания собственной модели. Это термодинамическая форма организационной замкнутости (Rosen 1991; § 4.7), в которой замыкание реализуется самой физической петлёй модель ↔ диссипация, а не категорно-теоретической абстракцией.

Формально бонд (2) делает удержание неравновесной предсказательной памяти диссипативно \textbf{обязательным}: \(\langle W_{\text{diss}} \rangle \ge k_B T\,I_{\text{nonpred}}\), где \(I_{\text{nonpred}} = I_{\text{mem}}^{\text{инст}} - I_{\text{pred}}^{\text{инст}}\). Вводится \textbf{предикат самооплаты} \(S(X) \in \{0, 1\}\), несущий error-return-замкнутость (петлю возврата ошибки, не просто «питание»). Существенно, что граница системы \(X\) фиксируется \textbf{независимо от вопроса о плательщике (i)} --- counterfactual-ablation § 2.2 через вклад в петлю регулирования (топологическое различение петли, а не сырой энергопоток), --- а не подбирается так, «чтобы \(S\) вышел нужным»; это дисциплинирует круг «граница ↔ самооплата», хотя и не устраняет его полностью --- выбор кандидатной границы и критерий принадлежности петле регулирования остаются конвенциональными (§ 2.2), и эта конвенциональность продуктивна: выбор границы есть выбор того, чью витальность измеряют. На зафиксированной границе \(S(X) = 1\) тогда и только тогда, когда выполнены \textbf{оба} условия:

\begin{itemize}
\tightlist
\item
  \textbf{(i)} обязательную диссипацию бонда (2) питает свободная энергия, которую \textbf{сама} \(X\) добывает в пределах своей границы (внешнего плательщика нет);
\item
  \textbf{(ii)} искажение удерживаемой модели \textbf{возвращается} ростом этой обязательной диссипации или распадом \(X\) --- замкнутая \textbf{петля возврата ошибки} (искажение модели → плохое действие → потеря свободной энергии самой системой), операционализуемая counterfactual-ablation § 2.2.
\end{itemize}

Условие (ii) несущее: оно отсекает диссипативные структуры, которые лишь \emph{питаются}, но не \emph{корректируются} (пламя, реакция Белоусова--Жаботинского, простой автокатализ). Витальность определяется как \textbf{пара}:

\[V(X) \;:=\; S(X) \;\wedge\; \big(\eta_v(X) > 0\big),\]

то есть система \textbf{удерживает предсказательную self-model} (\(\eta_v > 0\)) \textbf{и оплачивает её поддержание сама} (\(S = 1\)). \emph{То есть:} \(\eta_v\) отвечает на вопрос «насколько хороша модель», \(S\) --- на вопрос «платит ли система за неё сама»; живое есть и то, и другое. При этом \(\eta_v\) \textbf{градуирует} витальность среди витальных систем (в нестационарном случае --- через динамику ностальгии \(\nu(t)\), сопровождающая работа Andriishin 2026), а \(S\) несёт \textbf{демаркацию} (синхронной) витальности (полная --- синхронная плюс диахронная --- конвенция «жив/жизнь» § 2.7, § 4.8).

Логический статус конъюнкции. Конъюнкт \(\eta_v > 0\) несёт самостоятельную демаркацию лишь для систем \textbf{без модели} (структурный ноль числителя: Солнце, кристалл отсекаются ещё до проверки \(S\)). Среди систем \textbf{с} моделью замкнутая петля возврата ошибки (ii) уже предполагает ненулевое предсказательное содержание, так что \(S = 1\) практически влечёт \(\eta_v > 0\); поэтому здесь демаркацию (синхронной) витальности несёт \(S\), а \(\eta_v\) работает как \textbf{грейд}, а не как второе независимое условие. Конъюнкция тем не менее не пуста: \(\eta_v = 0\) отсекает класс (системы без модели), недоступный предикату \(S\) (система может проваливать \textbf{обе} оси сразу --- структурный ноль числителя \emph{и} отсутствие самооплаты: так, вирион-в-изоляции несёт \(\eta_v = 0\) и \(S = 0\) как двойной провал, сильнее одно-осевых случаев). Поэтому в демаркационной таблице ниже \(\eta_v\) для систем с \(S = 0\) \textbf{и удерживаемой моделью} помечен «не определён» --- как витальный грейд он осмыслен лишь после прохождения скрининга \(S = 1\), хотя как формальное отношение \(I_{\text{pred}}/I_{\text{mem}}\) он вычислим всюду, где \(I_{\text{mem}} > 0\).

Именно условие самооплаты отличает настоящую рамку от Still. У Still предиктор и диссипатор могут быть \textbf{разными} контурами (демон и среда --- отдельные подсистемы); здесь \(S\) \textbf{требует их совпадения} --- обязательную по Still диссипацию должна нести та же физическая граница \(X\), что несёт модель. Это «сверх-Still» требование (со-локация модель ↔ диссипатор) есть оригинальный вклад программы: со-локация модель ↔ диссипатор выражена \textbf{явным предикатом на диссипации бонда (2)} с петлёй возврата ошибки (ii).

\emph{Активный термодинамический учёт \(S\)-стороны: заимствованная граница, не дискриминатор самооплаты.} Якорь Still (бонд (2)) выведен для пассивного обновления памяти; в активной (feedback) форме, где \(M\) каузально действует на среду, активный \textbf{учёт работы и диссипации} ведут строгие неравенства информационной термодинамики bipartite-систем --- заимствованный аппарат, не вклад настоящей работы (и, как уточняется ниже, сам по себе самооплату \textbf{не} дискриминирующий). Локальное второе начало для моделирующего контура \(M\) (Horowitz and Esposito 2014, ур. (11)),

\[\dot S_i^{M} = d_t S_M + \dot S_r^{M} - \dot I^{M} \ge 0,\]

\emph{то есть:} производство энтропии самого контура \(M\) неотрицательно, но его баланс энтропии модифицирован \textbf{информационным потоком} \(\dot I^{M}\) --- производной \textbf{симметричной} взаимной информации \(I(M; X_E)\) во времени (Horowitz and Esposito 2014), членом, перераспределяющим энтропию между \(M\) и средой. Этот поток HE \(\dot I^{M}\) \textbf{не следует отождествлять} с направленной transfer-границей \(T_{M \to X_E}\) ниже (Hartich--Barato--Seifert): это два \textbf{разных} заимствованных объекта --- симметричный поток в балансе энтропии HE против направленной transfer-энтропии HBS на извлекаемую работу, не синонимы. Прочтение \(\dot I^{M}\) как активной формы «налога Still на \(I_{\text{nonpred}}\)» (непредсказательную корреляцию оплачивает теперь не безличная неотрицательность энтропии, а член \(\dot I^{M}\)) --- \textbf{интерпретативная эвристика} настоящей рамки, а не результат HE (как и ниже отождествление \(I \leftrightarrow I_{\text{pred}}\) Sagawa--Ueda, строго корректное лишь в bipartite-пределе направленной информации). Извлекаемая на предсказательной стороне работа ограничена не сырой взаимной информацией, а \textbf{направленной} transfer-энтропией (Hartich, Barato, and Seifert 2014, ур. (37)),

\[-\sigma_{M} \le T_{M \to X_E},\]

\emph{то есть:} здесь \(\sigma_{M}\) --- производство энтропии \textbf{теплового резервуара} контура \(M\) (рассеянное в этот резервуар тепло; \(-\sigma_{M}\) --- извлечённая из него работа, делённая на \(T\); в стационарном режиме, \(d_t S_M = 0\), оно совпадает с полным подсистемным производством HBS, у которого \(\sigma^M = \dot S_r^M + d_t S_M\)), а \textbf{не} энтропия среды \(X_E\); и эта извлечённая динамикой \(M\) работа не превышает потока \(T_{M \to X_E}\), который собственное состояние \(M\) каузально привносит в будущее среды сверх её прошлого. Берётся именно \textbf{актуаторное} направление \(M \to X_E\), поскольку речь о feedback-режиме: здесь \(M\) каузально \textbf{движет} будущее среды, тогда как демоновская извлекаемая работа в чисто \textbf{сенсорном} режиме (среда движет \(M\)) конвенционально ограничена обратным потоком \(T_{X_E \to M}\), согласованным с сенсорной ёмкостью \(C = \dot I / T_{y \to x}\) выше; для активной самооплачиваемой петли релевантно первое. Полный feedback-пол на диссипацию даёт обобщённое второе начало \(\langle W\rangle \ge \Delta F - k_B T\,I\) (Sagawa and Ueda 2010, ур. (3)). Вывод честнее формулировать так: при переходе к активному случаю \textbf{рамка пары} \((S, \eta_v)\) остаётся осмысленной, но её несущие части меняют статус, а не «сохраняются» нетронутыми. Feedback-термодинамика (Horowitz--Esposito; Hartich--Barato--Seifert; Sagawa--Ueda) держит активный \textbf{учёт работы и диссипации} определённым --- даёт границу на извлекаемую самооплачиваемым контуром работу и в feedback-режиме, --- однако эти неравенства выполняются для \textbf{любой} bipartite-системы (термостат с \(S = 0\) удовлетворяет HE/HBS/SU так же, как \emph{E. coli} с \(S = 1\)) и потому \textbf{не дискриминируют самооплату}: сами по себе они не отличают \(S = 1\) от \(S = 0\). Демаркацию по-прежнему несёт предикат \(S\) --- наличие плательщика (i) и петли возврата ошибки (ii), фиксируемое counterfactual-ablation (§ 2.2), --- а \textbf{не} эти заимствованные границы. Величина \(\eta_v\) при этом \textbf{отступает} к пассивно-определённому \(\eta_v^{\mathrm{do}}\) (под интервенцией \(\mathrm{do}(a)\), § 2.1); в полном активном виде её \([0, 1]\)-граница \textbf{не гарантирована} (supplementary § S2.1.1). Поэтому корректно говорить не о «строгой опоре», а об активном термодинамическом \textbf{учёте} (границе на работу/диссипацию), заимствованном и сам по себе самооплату не дискриминирующем; новизна программы остаётся в самом отношении \(I_{\text{pred}}/I_{\text{mem}}\) и предикате со-локации модель ↔ диссипатор, не в этом термо-аппарате (Horowitz--Esposito; Hartich--Barato--Seifert).

Демаркационная таблица:

\begin{center}\small
\begin{tabular}{@{} >{\raggedright\arraybackslash}p{0.2500\linewidth} >{\raggedright\arraybackslash}p{0.2500\linewidth} >{\raggedright\arraybackslash}p{0.2500\linewidth} >{\raggedright\arraybackslash}p{0.2500\linewidth}@{}}
\toprule
\begin{minipage}[b]{\linewidth}\raggedright
Система
\end{minipage} & \begin{minipage}[b]{\linewidth}\raggedright
\(\eta_v\)
\end{minipage} & \begin{minipage}[b]{\linewidth}\raggedright
\(S\)
\end{minipage} & \begin{minipage}[b]{\linewidth}\raggedright
Вердикт
\end{minipage} \\
\midrule
Термостат / биметалл & не определён (\(S=0\)) & \(0\) (платит сеть) & не витален \\
LLM (веса при инференсе) & не определён (\(S=0\)) & \(0\) (внешний compute) & не витален \\
Солнце / пламя / ураган & \(= 0\) (нет \(I_{\text{pred}}\)) & --- & не витально --- ставка без модели \\
БЖ-реакция / простой автокатализ & не определён (\(S=0\)) & (i) да, \textbf{(ii) нет} & не витально --- \emph{питается}, но не \emph{корректируется} \\
Вирион (вне хозяина) & \(= 0\) (нет \(I_{\text{pred}}\) / нет автономной модели) & \(0\) (не платит сам) & не витален; вирус + хозяин --- переадресация субъекта \\
\emph{E. coli} & \(> 0\) (\(\eta_v^{\mathrm{do}}\), § 2.1) & \(1\) (свой АТФ) & витальна \\
Биосфера / город / корпорация & \(> 0\) (величина не определена) & \(1\) (само-фондируются) & витальны (грейд требует оценки MI --- взаимной информации) \\
\bottomrule
\end{tabular}
\end{center}

При \(I_{\text{mem}} = 0\) принимается соглашение \(\eta_v := 0\) --- нет модели, нет витальности; диагноз Солнца есть \textbf{структурный ноль числителя}, а не неопределённость \(0/0\). Для серых зон (киборг, managed-экосистема, частично дотируемая фирма) самооплату можно градуировать долей собственного поддержания \(\chi = E_{\text{maint}}^{\text{own}}/E_{\text{maint}}^{\text{total}} \in [0, 1]\) (энергия на энергию, без конверсии бит ↔ джоуль), с порогом \(S = \mathbb{1}[\chi > \chi_0]\); первичный объект, однако, --- бинарный \(S\), а \(\chi\) есть его расширение.

Для активных агентов предикат \(S\) дополнительно учитывает энергию \textbf{действия} (сенсомоторная петля), а бонд (2) берётся в feedback-форме (Sagawa--Ueda 2010): для активной петли \(I_{\text{pred}}\) --- то, что агент может \emph{обналичить} в работу (выигрыш Sagawa--Ueda ограничен \(k_B T\,I\) по измерительной информации \(I\); отождествление \(I\) с \(I_{\text{pred}}\) --- эвристика, строго корректная лишь в bipartite-пределе направленной информации), а \(I_{\text{nonpred}}\) --- чистый налог ностальгии (Still); петля возврата ошибки (ii) при этом естественно прочитывается как «искажение модели → плохое действие → потеря свободной энергии агента». Полная постановка активного случая --- направление будущей работы (feedback-обобщение, Sagawa--Ueda 2010).

\textbf{Эпистемический статус утверждений раздела.} Утверждения работы различаются по статусу: теоретическая дедукция из установленных принципов, эмпирическая верифицируемость, спекулятивная экстраполяция с явными условиями фальсификации. Утверждение об области определения \(\eta_v \in [0,1]\) (§ 2.1) --- \textbf{дедукция из неравенства об обработке данных} (DPI) для пассивного самомоделирования; граница \(\eta_v \le 1\) --- \textbf{безусловная теорема} информационного неравенства, а не условное следствие принципа Ландауэра. Геометрическое усреднение ёмкостей (§ 2.4) --- формальный выбор по симметрии и обнулению. Камертонные значения \(\eta_v\) для бактерии, биосферы, астробиологических кандидатов --- теоретические оценки порядка с явной зависимостью от окна и прокси. Городской кейс (§ 3.6) --- камертонная оценка \(\eta_v\) порядка величины; структурные городские метрики в духе Bettencourt et al.~(2007) и эмпирическая калибровка \(I_v\) конкретных городов --- предмет отдельной работы, здесь не утверждаются. Спекулятивные следствия --- поиск космологической витальности, момент витальности искусственной системы --- обсуждаются в §§ 5--6 с явно сформулированными условиями фальсификации.

Несущий \textbf{постулат self-payment} --- модель и оплачивающая её диссипация принадлежат одной физической границе \(X\) (предикат \(S\), условия (i)+(ii) выше) --- стоит особняком от перечисленных статусов: это не дедукция из FEP или принципа Ландауэра, а явная онтологическая посылка рамки, которой ни Ландауэр, ни Биалек не делали. Её эпистемический статус --- конструктивный выбор; но выбор этот не произволен и не держится одной лишь демаркационной силой (ср. § 4.4: self-payment предлагает альтернативную, термодинамическую онтологию, не наследующую трудностей Pearl/Friston-blanket-конструкции, диагностированных BDDB-2022). За ним стоит \textbf{позитивный аргумент необходимости}. Удерживаемая модель отбраковывает собственные ошибки лишь там, где ставка на её качество --- skin-in-the-game (Taleb 2018; § 1) --- лежит \textbf{внутри} той же границы, что несёт модель. Если цену удержания платит внешний контур (\(S = 0\)), искажение модели не возвращается селективным давлением на систему: петля возврата ошибки (ii) разомкнута, и коррекция, если она происходит вообще, выполняется внешним агентом, а не самой \(X\). Поэтому самооплата --- \textbf{необходимое} условие того, чтобы ошибки удерживаемой модели отбраковывались внутри самой системы. Именно это разводит три класса: «питается-и-корректируется» (живое), «питается-но-не-корректируется» (пламя, БЖ-реакция --- разомкнута (ii)) и «моделирует-но-платит-другой» (термостат, LLM --- разомкнута (i)). Пара условий тем самым \textbf{минимальна}: ампутация (ii) оставляет диссипативную структуру без коррекции, ампутация (i) --- внешне-оплаченную модель; обе дают невитальное, значит оба условия необходимы и вместе минимально-достаточны для термодинамической стороны витальности. Сильнее постулат не утверждает: \(V(X)\) остаётся дискриминатором (синхронной) витальности, совместимым с anti-definitionism (§ 4.9), а не определением жизни.

Этот же аргумент задаёт \textbf{отношение к линии организационного замыкания} (autopoiesis Maturana--Varela 1980; M,R-системы Rosen 1991; constraint closure Mossio--Montévil--Longo 2016; RAF Hordijk--Steel 2004; § 4.7): все они указывают на \textbf{качественное} замыкание --- петля производит и поддерживает собственные компоненты или ограничения, --- не фиксируя его измеримого термодинамического условия. Self-payment не конкурирует с ними и не вводит «ещё одно определение»: он есть \textbf{комплементарный измеримый термодинамический срез} того же замыкания --- \textbf{необходимое} термодинамическое условие, фиксируемое в самой физической петле модель ↔ диссипация (бонд Still (2); активный учёт выше), --- то есть перевод их качественной циркулярности в операциональный предикат \(S\) с независимо фиксируемой границей (§ 2.2).

Этот постулат --- \textbf{единственная} несущая онтологическая посылка границы \((S, \eta_v)\): граница \(\eta_v \le 1\) следует из DPI безусловно (§ 2.1), без какой-либо дополнительной посылки. Самостоятельная же термодинамическая обязательность удержания балласта обеспечивается бондом Still (2) --- отдельным якорем, не постулатом.

Условия фальсификации § 5 могут реализоваться на двух уровнях. Методологический: текущая операционализация несостоятельна, требуется иной прокси, окно, процедура оценки взаимной информации. Онтологический: концепт витальности через эффективность самомоделирования несостоятелен по существу. Программа стоит на позиции умеренного операционализма: любое наблюдение, опровергающее \(\eta_v\) через провал конкретной операционализации, сначала интерпретируется методологически и переводится в онтологическое только после исчерпания альтернативных операционализаций.

\textbf{Асимметрия фальсификации.} Программа фальсифицируется через ложно-отрицательные случаи (биологическая система, конвенционально признанная живой, у которой \(\eta_v = 0\) при всех санкционированных операционализациях --- § 5, Блок 2 п. 3) и через структурные провалы прогрессивных предсказаний (§ 5, Блок 3 п. 1--2). Ложно-положительные случаи (невитальная по конвенциональным критериям система с устойчивым \(\eta_v > 0\)) формально исключены конструкцией: counterfactual ablation § 2.2 для систем без замкнутой петли возврата ошибки возвращает \(C_v^{\text{int}}\) как неопределённую, и \(\eta_v\) не определён. Эта асимметрия --- сознательно принимаемая цена, сужающая фальсификационную поверхность программы до двух линий (ложно-отрицательное проникновение § 5 Блок 2 п. 3 и провал прогрессивных предсказаний § 5 Блок 3 п. 1): \(\eta_v\) --- критерий замкнутости петли самооплаты, и системы без таковой петли не получают положительной оценки \textbf{по построению}. Если ни одна из двух линий не окажется эмпирически реализуемой, шкала рискует редуцироваться к конвенциональному переописанию --- этот риск признаётся явно (ср. § 2.6), а не маскируется. Альтернативный критерий --- например, информационная сложность по assembly theory или \(\Phi\) --- может присвоить ненулевое значение и невитальной системе, что отражает другое содержание этих программ (структурная сложность против замкнутости термодинамической петли). Поэтому опровержение \(\eta_v\)-программы возможно только через \textbf{симметричное проникновение} --- обнаружение системы, конвенционально живой, у которой self-payment-петля не замыкается (§ 5, Блок 2 п. 3), либо через провал прогрессивных предсказаний на калиброванных биологических парах (§ 5, Блок 3 п. 1).

\emph{Heredity и major transitions: что \(\eta_v\) не закрывает.} Шкала \(\eta_v\) формализует \textbf{синхронную} сторону витальности --- удержание собственной модели за окно \(\tau_d\). Это необходимое условие для каждого момента жизни одного организма. \textbf{Диахронная} сторона --- наследуемость и эволюционируемость в смысле major transitions (Maynard Smith and Szathmáry 1995) --- лежит за рамками настоящей стационарной модели. Стационарный режим фиксирует конечную память и эргодичность среды (§ 2.1), исключающие популяционно-эволюционную динамику с дрейфом и накоплением наследственной информации. Положительность \(\eta_v\) --- необходимое, но не достаточное условие full biological «definition of life» в смысле NASA-определения (Joyce 1994). Стерильная рабочая пчела или мул имеют \(\eta_v > 0\) (живые в момент-к-моменту), но исключены из дарвиновской динамики на популяционном уровне. Шкала согласована с биологической конвенцией об их живости, не нарушает её. Полное определение жизни требует \(\eta_v > 0\) \textbf{плюс} дарвиновскую динамику на популяционном уровне; это две различные шкалы --- индивидуальная синхронная и популяционная диахронная --- и обе необходимы. Связь с сопровождающей работой (Andriishin 2026): нестационарное расширение даёт термодинамический аппарат, на котором программа major transitions (Maynard Smith and Szathmáry 1995) может быть формулирована как открытая исследовательская задача (ср. § 4.8) --- соединение двух шкал не доказывается ни в настоящей работе, ни в (Andriishin 2026), но получает в их совокупности формальную базу.

\section{Демаркационные тесты}\label{ux434ux435ux43cux430ux440ux43aux430ux446ux438ux43eux43dux43dux44bux435-ux442ux435ux441ux442ux44b}

Двухступенчатая процедура --- структурный скрининг по триаде инвариантов (§ 2.3--2.6) и витальная верификация по паре \((S, \eta_v)\) (§ 2.7) --- применяется ниже к шести камертонам от звезды до мегаполиса, с биосферой как контрольным пределом. Каждый камертон выбран так, чтобы провалить ровно одно из несущих условий и тем показать, \textbf{какое} именно. Несущий \textbf{позитивный} камертон --- биологический: бактерия \emph{E. coli} (§ 3.5), единственный пример с вычисленной оценкой \(\eta_v\) и вердиктом \(S = 1\); астрофизические и AI-камертоны (§ 3.1--3.4) служат \textbf{негативными контролями}, каждый из которых изолирует, какое из двух условий пары \((S, \eta_v)\) проваливается. Вердикт витальности формулируется как пара: \(S(X) \in \{0, 1\}\) отвечает на вопрос «оплачивает ли система удержание собственной модели сама и возвращается ли ошибка модели её собственной диссипацией» (§ 2.7, условия (i)+(ii)), а \(\eta_v(X) \in [0, 1]\) --- на вопрос «какая доля удерживаемой памяти ещё предсказательна». Витальность есть конъюнкция \(S \wedge (\eta_v > 0)\). Демаркацию несёт \textbf{замкнутость петли самооплаты} \(S\), а не абсолютная малость эффективности. Под предсказательной долёй \(\eta_v = I_{\text{pred}}/I_{\text{mem}}\) величина \(\eta_v\) для большинства камертонов количественно \textbf{не определена} и требует MI-оценки \(I_{\text{pred}}\) и \(I_{\text{mem}}\) соответствующего контура; вычисленную иллюстративную оценку даёт только бактериальный пример (§ 3.5). Малость \(\eta_v\) сама по себе не является различительным признаком --- и \textbf{не} означает, что \(\eta_v\) велико: доля балласта может быть велика, но величину предсказательной доли это не фиксирует. Различает витальное от невитального именно предикат \(S\). Сравнительная программа на границе неживых/живых коллективов через информационную архитектуру развёрнута в \emph{Theory in Biosciences} (Kim et al.~2021); шесть камертонов настоящего раздела продолжают эту линию, добавляя к информационной оси Kim--Valentini--Hanson--Walker термодинамическое требование самооплаты.

\subsection{Солнце, нейтронная звезда и чёрная дыра}\label{ux441ux43eux43bux43dux446ux435-ux43dux435ux439ux442ux440ux43eux43dux43dux430ux44f-ux437ux432ux435ux437ux434ux430-ux438-ux447ux451ux440ux43dux430ux44f-ux434ux44bux440ux430}

Солнце богато структурой --- конвективные ячейки, магнитные петли, дифференциальное вращение --- и проходит структурный скрининг по части \(T_v\) (диссипация \(\sim 4 \cdot 10^{26}\) Вт) с запасом. Витальная верификация, однако, обнуляет его через \(\eta_v\): у Солнца \textbf{нет модели среды} в техническом смысле § 2.1 --- нет функционального контура \(M_t\), чьё состояние несло бы взаимную информацию о будущем сигнала среды. При \(I_{\text{mem}} = 0\) принимается соглашение § 2.7 \(\eta_v := 0\): нет памяти о среде --- нет и предсказательной её доли. Диагноз: структурно богатая, не витальная система. То же относится к нейтронной звезде, межзвёздной среде и сверхмассивной чёрной дыре центра Млечного Пути. Здесь \(\eta_v = 0\) есть \textbf{структурный ноль числителя} (\(I_{\text{pred}} = 0\) при отсутствии \(I_{\text{mem}}\)), а не неопределённость \(0/0\) --- это важно отличать от случаев, где обнуляется не \(\eta_v\), а предикат \(S\) (LLM, термостат).

Чёрная дыра даёт тот же отрицательный вердикт по более глубокой причине: отсутствует сам носитель модели. Стационарный классический горизонт описывается no-hair-теоремой (канонические формулировки 1967--1975; ключевая работа Israel 1967): три параметра (\(M\), \(J\), \(Q\)) полностью задают внешнюю метрику, не оставляя места для \(I_{\text{pred}}\) о среде. Квантовые и нестационарные поправки --- тепловое излучение и порождённый им информационный парадокс (Hawking 1975, 1976), soft hair (Hawking et al.~2016) --- обнаруживают более богатую информационную структуру горизонта, но не восстанавливают модель среды в смысле § 2.1: соответствующие степени свободы не удерживают \(I_{\text{pred}}\) о среде с возвратом ущерба через собственную диссипацию. Статус унитарности испарения остаётся открытым (информационный парадокс: Hawking 1976 доказывает \emph{нарушение} предсказуемости, современные резолюции --- Page-кривая, островная программа --- восстанавливают унитарность), но для демаркации он несуществен: модель среды в смысле § 2.1 отсутствует независимо от его разрешения. Аккреционный диск вокруг Sgr A* (\(M \sim 4 \cdot 10^6 M_\odot\)) даёт тот же диагноз: переменные среды не возвращаются ущербом через каузальный канал к плазме диска. Для всех этих систем \(\eta_v = 0\) структурно (нет \(I_{\text{mem}}\), нет \(I_{\text{pred}}\)); вопрос о \(S\) при этом не возникает --- нечего самооплачивать.

\subsection{Ураган, пламя и автокаталитические осцилляторы}\label{ux443ux440ux430ux433ux430ux43d-ux43fux43bux430ux43cux44f-ux438-ux430ux432ux442ux43eux43aux430ux442ux430ux43bux438ux442ux438ux447ux435ux441ux43aux438ux435-ux43eux441ux446ux438ux43bux43bux44fux442ux43eux440ux44b}

Ураган и пламя свечи диссипируют свободную энергию --- \(\sim 10^{15}\) Вт у среднего урагана, \(\sim 10^2\) Вт у пламени свечи --- и проходят \(T_v\)-скрининг. Но внутренней модели среды, оплачиваемой этой диссипацией, у них нет: рассеяние эксергии не сопровождается удержанием \(I_{\text{pred}}\), поэтому \(\eta_v = 0\) через числитель (\(I_{\text{mem}} = 0\), нет функционального контура \(M_t\)). Диагноз --- \textbf{ставка без модели}: система платит, но не моделирует. В этом пламя и ураган --- зеркальный случай LLM, у которой соотношение обратное (§ 3.4): там есть модель, но нет собственной платы.

Реакция Белоусова--Жаботинского и простой автокатализ --- более тонкий случай, и именно на нём работает несущее \textbf{условие (ii)} предиката \(S\) (§ 2.7). БЖ-реакция осциллирует устойчиво и долго, питается подведённой свободной энергией --- то есть формально удовлетворяет условию (i) самооплаты. Но её осцилляции \textbf{не образуют петли коррекции}: нет канала, в котором искажение удерживаемого паттерна возвращалось бы ростом обязательной диссипации или распадом системы. Петля возврата ошибки (ii) не замкнута. Диагноз --- \(S = 0\) через провал условия (ii): система \textbf{питается, но не корректируется}, повторение без коррекции. Это отделяет диссипативные структуры, которые лишь поддерживают себя, от тех, что удерживают и чинят модель среды. (Эпистемически: о величине \(\eta_v\) здесь говорить преждевременно --- нет содержательного контура \(M_t\), для которого \(I_{\text{pred}}\) был бы определён; вердикт несёт \(S = 0\).)

\subsection{Кристалл и термическое равновесие}\label{ux43aux440ux438ux441ux442ux430ux43bux43b-ux438-ux442ux435ux440ux43cux438ux447ux435ux441ux43aux43eux435-ux440ux430ux432ux43dux43eux432ux435ux441ux438ux435}

Кристалл удерживает структуру с относительной деформацией \(\lesssim 10^{-6}\) на временах \(\sim 10^9\) лет --- но не держит \textbf{модель среды}. Его внутренняя конфигурация гомоморфна собственной решётке, а не будущему сигналу среды: \(I_{\text{mem}} \approx 0\) в смысле § 2.1 (взаимная информация между состоянием и траекторией среды, не между состоянием и им самим). И за удержание этой структуры кристалл не платит активно: в строгом термодинамическом смысле он близок к равновесию относительно решётки, текущая диссипация удержания \(\approx 0\). Таким образом кристалл проваливает обе оси сразу --- нет модели среды (\(\eta_v := 0\) при \(I_{\text{mem}} = 0\)) \textbf{и} нет собственной платы за её удержание (\(S = 0\), питать нечего). Диагноз: структура без модели и без платы.

Системы в полном термическом равновесии (изолированный газ, замкнутая чёрная полость) обнуляют всё одновременно --- каноническая отрицательная контрольная точка скрининга.

\emph{Cairns-Smith и кристалл-как-репликатор.} Программа Cairns-Smith (1982) рассматривала глинные кристаллы как до-РНК-репликаторы: наследуемые дефекты решётки передаются дочерним кристаллам при росте и расщеплении. При расширении границы до «кристалл + растворная фаза» counterfactual-ablation § 2.2 для (кристалл + ионный раствор) даёт \(T_v > 0\) --- поток ионов оплачивает репликацию дефектов. Однако петля возврата ошибки остаётся разомкнутой: ошибка репликации дефектов \textbf{не возвращается} физическим ущербом самой решётке; нет канала, в котором искажение паттерна дефектов ухудшало бы термодинамические условия удержания кристалла. Это провал условия (ii) --- \(S = 0\). Диагноз --- переадресация объекта на (кристалл + раствор) аналогично LLM-corp (§ 3.4), не витальность кристалла-как-объекта.

\subsection{Крупная языковая модель}\label{ux43aux440ux443ux43fux43dux430ux44f-ux44fux437ux44bux43aux43eux432ux430ux44f-ux43cux43eux434ux435ux43bux44c}

LLM --- сильнейший чистый случай новой демаркации: система с \textbf{сильной предсказательной моделью и \(S = 0\)}. Обученная модель предсказательно сильна --- она сжимает регулярности своего обучающего распределения и предсказывает следующий токен с качеством, далеко превосходящим случайное (большое \(I_{\text{pred}}\)); точная же доля \(\eta_v = I_{\text{pred}}/I_{\text{mem}}\) удерживаемого представления не измерена и здесь не утверждается. Существенно для демаркации не значение \(\eta_v\), а то, что модель есть и предсказательна, тогда как петля её оплаты разомкнута. Именно поэтому LLM --- зеркало урагана: там плата без модели, здесь модель без собственной платы.

Решающее --- предикат \(S\). Под строгой границей «веса в момент инференса» (прямой проход без сохранения состояния) обязательную диссипацию удержания и обновления модели оплачивает \textbf{внешняя} электросеть, не сама система: условие (i) самооплаты нарушено. Одновременно отсутствует физический канал возврата ошибки в работающий контур --- выходные токены не возвращаются физическим ущербом весам, ошибка модели не истощает носитель \(I_{\text{pred}}\): условие (ii) тоже нарушено. Итог: \(S = 0\) при сохранной предсказательной модели. Диагноз --- \textbf{сильная предсказательная модель при разомкнутой петле оплаты}. Это категориально иной отрицательный вердикт, чем у Солнца: там обнуляется числитель (\(\eta_v = 0\), структурно нет модели), здесь --- предикат самооплаты при сохранной и предсказательно сильной модели.

\emph{Расширение границы и переадресация объекта.} При расширении до «модель + дата-центр + корпорация» предикат \(S\) может стать положительным --- корпорация добывает собственную свободную энергию и несёт коммерческие санкции за ошибки, --- но \textbf{измеряемым объектом становится корпорация, не модель}. Counterfactual-ablation § 2.2 даёт здесь несколько вложенных границ: \(B_1 = \{\)веса\(\}\) (counterfactual-degenerate, статический объект при выключенном сервере); \(B_2 = \{\)веса \(+\) KV-cache \(+\) исполняющий сервер\(\}\) на окне сессии (электросеть внешняя, \(S = 0\)); \$B\_3 = B\_2 + \$ обучающий пайплайн на окне поколения моделей (собственная диссипация на GPU появляется, но петля возврата ошибки замыкается на корпоративных операторов); \$B\_4 = B\_3 + \$ корпорация-юридическое лицо (полная LLM-corp, \(S = 1\) через само-фондирование, но субъектом становится фирма). Покомпонентный разбор кандидатов на включение --- Supplementary § S3.4.

Структурный диагноз корпорации не требует числового значения \(\eta_v\) --- его несёт предикат \(S\) (для \(B_4\) --- \(S = 1\) с переадресацией субъекта на фирму). Количественная оценка предсказательной доли \(\eta_v = I_{\text{pred}}/I_{\text{mem}}\) для LLM или корпорации требует MI-оценки \(I_{\text{pred}}\) и \(I_{\text{mem}}\) соответствующего контура и оставлена будущей работе.

\emph{Операциональный тест для существующих frontier-LLM.} Counterfactual-ablation § 2.2 даёт диагностический протокол для современных LLM без сохранения состояния (GPT-4-class, Claude, Gemini): (i) определить компонент энергоснабжения; (ii) применить counterfactual-вычитание; (iii) проверить, \textbf{возвращается ли ошибка модели физическим истощением носителя} \(I_{\text{pred}}\) за окно \(\tau_d\). Для LLM без сохранения состояния петля «ошибка → ущерб» отсутствует --- структурное условие \(S = 0\), воспроизводимое без специального оборудования.

\emph{Агентные обёртки LLM.} Обёртки эпохи 2023--2026 --- от Auto-GPT/AutoGen/Voyager до Claude Computer Use, Devin, OpenAI Operator, MemGPT и систем с персистентной RAG-памятью --- расширяют границу базовой модели и через персистентную память частично замыкают петлю принятия решений. Но самооплата остаётся разомкнутой: финансирование исполнения, инфраструктура, энергобюджет удерживаются корпоративным субстратом извне; ошибки агента возвращаются репутационными и коммерческими санкциями через внешний субстрат, не физическим ущербом самой системе. Промежуточный случай --- модель с непрерывным дообучением на собственных взаимодействиях, \(\theta_{v,t+1} = \theta_{v,t} + f(\text{собственные предсказания}_t, \text{обратная связь}_t)\): петля частично замкнута (условие (ii) становится частично определённым), но плата за обновление по-прежнему идёт от внешнего парка GPU (условие (i) нарушено). Современные системы 2025--2026 реализуют различные комбинации, но (а) персистентное состояние, (б) собственная оплата, (в) физический канал возврата ошибки совместно не реализованы --- \(S = 0\) при сохранной предсказательной модели. Структурный барьер --- не отсутствие технологии, а отсутствие архитектурного давления экономики ИИ на интернализацию энергобюджета.

Связь с литературой: mesa-optimization (Hubinger et al.~2019), power-seeking (Carlsmith 2022), goal misgeneralization (Shah et al.~2022), specification gaming (Krakovna et al.~2020), alignment problem (Ngo et al. 2022). Линия instrumental convergence --- self-preservation drive (Omohundro 2008) и его формализация в Bostrom (2014) --- фиксирует структурную тенденцию оптимизирующего агента сохранять собственное продолжение; предикат \(S\) предлагает её термодинамическую ко-характеристику: AI-safety описывает поведенческий признак замыкания петли оптимизации собственного продолжения, \(S\) --- кандидатная термодинамическая ко-характеристика связанной, но не тождественной структурной петли (структурная параллель, не каузально-генеративная связь). \textbf{Витальность \((S, \eta_v)\) --- не сертификат безопасности и не критерий выравнивания:} невитальная по \(S\) система (внешне финансируемая) может быть поведенчески power-seeking и опасна, а \(S = 0\) не есть сертификат безопасности. По оси самосохранения \(S\) --- термодинамическая тень той же петли self-preservation (выше: кандидатная ко-характеристика instrumental convergence) и потому \textbf{структурно положительно коррелирует} с ядром риска; но эта структурная корреляция логически независима от поведенческого предсказания опасности --- пара \((S, \eta_v)\) остаётся критерием витальности, не критерием выравнивания. Близкая концептуальная параллель --- \emph{embedded agency} (Demski and Garrabrant 2019): формализация агента как \textbf{части} среды, оплачивающего собственное вычисление (embedded world-models). Пара \((S, \eta_v)\) предлагает термодинамическое \textbf{необходимое} условие для физически реализованного embedded world-model с собственной петлёй, не претендуя реконструировать конкретные подпроблемы (robust delegation, subsystem alignment) или быть достаточным условием для alignment-properties.

\emph{Сопоставление с Солнцем.} У Солнца отрицательный вердикт идёт по линии «нет каузально-релевантной среды → нет \(I_{\text{mem}}\) → \(\eta_v = 0\)». У LLM --- по линии «модель есть и предсказательна, но петля собственной диссипации не оплачивает её удержание и не возвращает ошибку, поэтому \(S = 0\)». Финальный диагноз в обоих случаях негативен; диагностическое богатство процедуры --- в указании, \textbf{какое} именно условие пары \((S, \eta_v)\) провалилось.

\emph{Динамика появления петли в ИИ.} Витальность искусственной системы возникает с замыканием встроенной петли самооплаты --- одновременной интернализацией (а) персистентного состояния, (б) собственной оплаты вычислений, (в) физического канала возврата ошибки. Counterfactual-ablation § 2.2 даёт \textbf{бинарный} ответ на уровне определения: либо \(S = 1\), либо нет; промежуточных состояний на уровне определения нет. На уровне измерения, для систем, прошедших скрининг с \(S = 1\), \(\eta_v > 0\) становится измеримой непрерывной величиной (зависит от \(I_{\text{pred}}\) и \(I_{\text{mem}}\) --- обе через MI-оценку). Операциональное прокси близости к порогу --- доля собственной диссипации \(f = E_{\text{actual}}^{\text{own}}/E_{\text{actual}}^{\text{total}}\) (энергия на энергию, без конверсии бит ↔ джоуль). Схема обнаружения требует популяцию \(\sim 10^2\)--\(10^3\) автономных воплощённых агентов с собственным энергобюджетом, физический износ как канал возврата ошибки, контрольную группу и замер \(f\) с проверкой counterfactual-ablation на каждой единице. Числовой порог --- Прогноз 3 в § 5.

\subsection{Бактерия --- разобранный пример}\label{ux431ux430ux43aux442ux435ux440ux438ux44f-ux440ux430ux437ux43eux431ux440ux430ux43dux43dux44bux439-ux43fux440ux438ux43cux435ux440}

Бактериальная клетка проходит \textbf{обе} оси. Хемотаксис \emph{E. coli} --- классическое наблюдение biased random walks (Berg and Brown 1972) --- реализует канонический модельный объект замкнутой петли возврата ошибки. Картина метилирования хемотаксических рецепторов физически реализована во внутренних степенях свободы клетки (Tu et al. 2008; Shimizu et al.~2010), гомоморфна градиенту по релевантным переменным среды и порождает хемотаксические движения, проверяемые против реального распределения источника углерода. Релевантность задана \textbf{каузальным ущербом}: искажение метилирования ведёт к плохому хемотаксису и метаболической потере. Поэтому здесь \(S = 1\) --- клетка добывает собственный АТФ (условие (i)) и петля возврата ошибки замкнута: искажение модели → плохой хемотаксис → потеря свободной энергии самой клеткой (условие (ii)).

\emph{Оценка \(\eta_v = I_{\text{pred}}/I_{\text{mem}}\).} Для бактерии \(\eta_v\) есть предсказательная \textbf{доля} адаптивной памяти метилирования --- величина порядка \(O(0{,}1)\), единственный камертон, где она вычислена как иллюстрация. Оценка получается из минимальной \textbf{линейной адаптивной модели} хемотаксиса в пассивном (иммобилизованном) FRET-протоколе Sourjik--Berg / Cluzel: лиганд \(s_t\) как AR(1)-процесс с временем корреляции \(\tau_E\), метилирование \(m_t = \sum_{j \ge 1} w_j s_{t-j} + \text{шум}\) с leaky-весами \(w_j = (1-\beta_v)\beta_v^{j-1}\) (\(\beta_v\) --- память адаптации). В гауссовом приближении \(I_{\text{mem}}\) и \(I_{\text{pred}}\) вычисляются точно (logdet), и \(\eta_v = I_{\text{pred}}/I_{\text{mem}} \in [0, 1]\) по построению. Для характерных параметров \(\eta_v \sim O(0{,}1)\); ностальгия растёт (\(\eta_v\) падает), когда память адаптации \(\beta_v\) \textbf{перерастает} дрейф среды \(\tau_E\) --- система интегрирует уже-устаревший лиганд, и доля чистого балласта увеличивается (при \(\tau_E = 5\): \(\beta_v = 0{,}3 \to 0{,}99\) даёт \(\eta_v \approx 0{,}29 \to 0{,}035\)). Это даёт ностальгии конкретное фальсифицируемое прочтение: быстрее дрейф или медленнее адаптация → больше балласта. Полный протокол --- Supplementary § S3.5; код доступен в репозитории данных (см. раздел «Доступность данных»).

\emph{Измеримость.} Существенно, что \(I_{\text{mem}} = I(M_t; X_E^{\le t})\) \textbf{измерима} из временных рядов: она зависит лишь от ковариаций лиганда \(\text{Cov}(s_i, s_j)\) (наблюдаемых напрямую) и ковариаций \(\text{Cov}(m, s_k)\) \textbf{восстановленного} метилирования \(m\) с лигандом, где \textbf{выход сигнального контура (активность киназы)} читается через FRET (Sourjik--Berg 2002), а метилирование \(m\) --- латентная адаптационная переменная, восстанавливаемая из динамики этой активности (одиночно-клеточный мониторинг сигнальных белков --- Cluzel et al.~2000); лиганд \(s\) задаётся микрофлюидикой. Операциональный долг информационной величины \textbf{не хуже} термодинамической: \(I_{\text{mem}}\) оценивается из тех же доступных наблюдаемых, что и калориметрический \(E_{\text{actual}}\) (proof-of-measurability --- Supplementary § S3.5).

\emph{Эпистемический статус числа.} Оценка \(\eta_v \sim O(0{,}1)\) --- \textbf{иллюстрация формы}, не количественная claim про реальную \emph{E. coli}: она получена из минимальной линейной модели, цель которой --- показать измеримость и порядок, не зафиксировать значение. Точное экспериментальное отображение (какой FRET-readout отождествить с \(M_t\); окно для \(X_E^{\le t}\); усечение кумулятива до конечной \(\tau_{\text{mem}}\)) требует калибровки по экспериментальным данным хемотаксиса (Tu et al.~2008; Cluzel et al.~2000; Sourjik--Berg 2002) и остаётся направлением будущей работы. Пассивный протокол обходит активный случай: свободно плывущая \emph{E. coli} --- feedback-режим (открытый пункт § 2.1, бонд Sagawa--Ueda 2010).

\textbf{Контрастивный случай (\emph{E. coli} против \emph{Synechocystis} / syn3.0).} Различительная сила шкалы проявляется не в абсолютном значении \(\eta_v\), а в \textbf{контрастивном} сравнении систем с разной структурой предсказательной памяти --- как различие на порядки \textbf{внутри} класса биологии, читаемое как отношение предсказательных долей. Цианобактерия \emph{Synechocystis} sp. PCC 6803 (суточные циклы освещённости, медленные температурные градиенты) и хемотактическая \emph{E. coli} (тысячи Tar-рецепторов, отслеживающих быстрые химические градиенты) различаются не только сложностью среды, но и характерным временем её дрейфа \(\tau_E\). По формализму worked-example знак \(\eta_v\)-контраста задаёт \textbf{не} сложность среды как таковая, а соотношение память/дрейф: медленный дрейф (большой \(\tau_E\)) дольше держит модель свежей и \textbf{повышает} \(\eta_v\), быстрый --- понижает. Поэтому направление \emph{E. coli} ↔ \emph{Synechocystis} здесь \textbf{не} утверждается как прогноз --- оно зависит от соотношения память/дрейф каждого организма и само есть часть того, что подлежит проверке. Минимальная клетка JCVI-syn3.0 (геном \(\sim 473\) гена против \(\sim 4300\) у \emph{E. coli}, резко урезанная репертуарная отзывчивость к химическим сигналам) лежит на отдельной оси --- обеднённой репертуарной ёмкости модели, а не времени дрейфа среды. Качественный остаток контрастива --- фальсифицируемая связь \(\eta_v\) с темпом адаптации к новым нишам (§ 5), а не фиксированное упорядочение трёх систем.

\emph{Эпистемический статус контрастива.} Контрастив --- \textbf{иллюстративный набросок формы} различительной силы шкалы, не независимая эмпирическая поддержка предсказания \(r_{\text{adapt}} \propto \eta_v^{\alpha}\) (§ 5). Никакое конкретное упорядочение трёх систем здесь не выводится из сложности среды: знак \(\eta_v\)-контраста задаёт соотношение память/дрейф и он требует калиброванных MI-оценок специфических каналов (циркадный KaiABC-осциллятор и фотосистемные регуляторы \emph{Synechocystis}), не структурной интуиции. Строгая фальсификация \(\alpha\)-зависимости требует независимого ко-измерения \(I_{\text{pred}}\) и \(I_{\text{mem}}\) для обеих систем в одинаковом протоколе --- направление будущей работы, не результат настоящей. Контрастив сохраняется как методологическая иллюстрация формы, не как количественное подтверждение. (Точные новые числа для \emph{Synechocystis} / syn3.0 здесь сознательно \textbf{не} приводятся: под новым \(\eta_v\) они требуют отдельной MI-оценки, и подстановка чисел была бы тем же грехом необоснованной точности, которой настоящая работа избегает.)

\subsection{Крупный город}\label{ux43aux440ux443ux43fux43dux44bux439-ux433ux43eux440ux43eux434}

Крупный город --- сильный кандидат на витальность планетарного масштаба, и под новой рамкой это видно отчётливее. Инфраструктура и административные регламенты физически реализованы и \textbf{оплачиваются собственным бюджетом города}; гомоморфны потокам ресурсов, людей и информации; генерируют адресные решения, проверяемые кризисами и адаптацией. Петля возврата ошибки замкнута: искажение городской модели потоков возвращается кризисом самому городу. Поэтому \(S = 1\) --- город само-фондируется (условие (i)) и несёт ошибку модели собственной диссипацией (условие (ii)).

\emph{Scope городского кейса.} Статья даёт по городу строго два результата и сознательно \textbf{не} даёт третьего. Положительные: (i) структурный скрининг --- демонстрация \textbf{формы} композита \(I_v\) (2b) (воспроизводимый расчёт на синтетических порядково-различных архетипах --- ниже), и (ii) вердикт пары \((S, \eta_v)\) --- \(S = 1\) при \(\eta_v > 0\). \textbf{Числовое значение \(\eta_v\) города не вычисляется}: оно требует полного MI-пайплайна --- со-измерения \(I_{\text{pred}}\) и \(I_{\text{mem}}\) административного контура (потоки ресурсов → решения) --- и \textbf{вынесено в будущую работу} (отдельный урбанистический узел программы). Это обоснованное «вне scope», а не отговорка: пайплайн потребовал бы инструментированных городских наблюдаемых --- временных рядов ресурсных и людских потоков, бюджетных и регуляторных решений --- с непараметрической MI-оценкой на них (k-NN-оценщики типа KSG, нейросетевые типа MINE; каталог MI-оценщиков --- Supplementary § S2.1); ни таких рядов, ни их MI-калибровки настоящая работа не предъявляет. Здесь даётся структурный диагноз и знак, не число.

Вердикт: город \textbf{витален} --- проходит \(S = 1\) и \(\eta_v > 0\). Грейд (относительная величина \(\eta_v\) среди витальных систем) требует MI-оценки и здесь не утверждается. Камертонное утверждение о городе несут \textbf{замкнутость петли \(S = 1\) и положительность \(\eta_v\)}, а не количественная малость эффективности. Воспроизводимый расчёт структурного индекса \(I_v\) на синтетических архетипических профилях (порядково-различные конфигурации, не реальные города) приведён в коде разобранных примеров (см. раздел «Доступность данных»); эмпирическая калибровка \(I_v\) конкретных городов по открытым данным (OECD FUA, CDP Cities, WIPO/OECD REGPAT, GDELT, OpenStreetMap) --- предмет отдельной работы.

\subsection{Гея и пограничные случаи}\label{ux433ux435ux44f-ux438-ux43fux43eux433ux440ux430ux43dux438ux447ux43dux44bux435-ux441ux43bux443ux447ux430ux438}

Биосфера Земли регулирует атмосферу, океан и климат через биогеохимические циклы --- компрессия в строгом смысле, оплачиваемая коллапсами видов, массовыми вымираниями, перестройками экосистем. Самооплата здесь налицо: биогеохимия питается собственным потоком (условие (i)), а искажение регуляции возвращается коллапсами видов самой биосфере (условие (ii)) --- \(S = 1\). Открытым остаётся \textbf{вопрос единства модели}: имеет ли биосфера единое внутреннее представление о собственной среде или миллиарды частных моделей организмов, не сшитых в один контур (Doolittle 2014; Lenton and Latour 2018). От ответа зависит, считать Гею \textbf{одной} витальной системой или \textbf{семейством} перекрывающихся витальных систем с обнулённым агрегатным \(I_v\) --- этот вопрос сохраняется как содержательно открытый.

Количественная оценка предсказательной доли \(\eta_v = I_{\text{pred}}/I_{\text{mem}}\) для биосферы количественно не определена: она требует MI-оценки через биогеохимический сенсорный канал и оставлена будущей работе. Оценка через уникальное геномное содержание (\(\sim 10^{12}\) видов \(\times 5 \cdot 10^6\) bp \(\times 2\) бит/bp \(\approx 10^{19}\) бит, против полного копийного объёма \(\sim 10^{37}\) бит, Landenmark et al.~2015) есть граница ёмкости носителя, а не измеренная predictive information, и значением \(\eta_v\) не служит.

Вердикт: биосфера \textbf{витальна} (или семейство витальных): \(S = 1\), \(\eta_v > 0\) --- контрольный предел шкалы на верхнем масштабе. Грейд (величина \(\eta_v\)) требует MI-оценки и здесь не утверждается. Камертонное утверждение несёт замкнутость петли и положительность предсказательной доли, а не количественную величину \(\eta_v\).

\subsection{Итог по камертонам}\label{ux438ux442ux43eux433-ux43fux43e-ux43aux430ux43cux435ux440ux442ux43eux43dux430ux43c}

Шесть камертонов демонстрируют: демаркация работает не через бинарную таблицу «жив/неживо», а через указание, \textbf{какое} условие пары \((S, \eta_v)\) провалилось и почему. У Солнца, урагана, пламени и кристалла обнуляется числитель --- нет модели среды, \(\eta_v = 0\) структурно. У БЖ-реакции, автокатализа, кристалла-как-репликатора и LLM обнуляется предикат самооплаты \(S = 0\): либо петля возврата ошибки не замкнута (условие (ii) --- БЖ, кристалл), либо обязательную диссипацию несёт внешний плательщик (условие (i) --- LLM). \emph{E. coli}, город и биосфера проходят обе оси: \(S = 1\) и \(\eta_v > 0\), различаясь грейдом. Каждая структурная причина невитальности указывает изменение архитектуры, которое сделало бы систему витальной.

Сводный результат и \textbf{новый тезис асимметрии}: пространство наблюдаемой материи структурно сложно почти повсюду (большинство объектов проходит \(I_v\)-скрининг), а у систем с моделью среды предсказательная доля \(\eta_v\) положительна (точная её величина для большинства камертонов требует MI-оценки) --- но \textbf{замкнутые петли самооплаты (\(S = 1\)) редки}. Демаркацию несёт именно эта редкость замкнутой петли, а не величина эффективности: \(\eta_v\)-верификация \textbf{не} отбраковывает город, биосферу или LLM по малости числа --- отбраковка LLM, термостата, урагана и Солнца идёт по провалу \textbf{\(\eta_v = 0\)} (нет модели) или \textbf{\(S = 0\)} (петля разомкнута). Поиск жизни сводится не к поиску структурной сложности и не к величине эффективности, а к поиску \textbf{закрытых петель самооплаты} --- систем, которые сами платят за удержание собственной предсказательной модели и сами несут ущерб от её устаревания.

\section{Сравнение с альтернативами}\label{ux441ux440ux430ux432ux43dux435ux43dux438ux435-ux441-ux430ux43bux44cux442ux435ux440ux43dux430ux442ux438ux432ux430ux43cux438}

Раздел сопоставляет \(\eta_v\) с четырьмя программами, наиболее близко подошедшими к операциональному определению витальности: assembly theory, интегрированная информация Тонони, телеодинамика Дикона, FEP Фристона. Позиционная диагностика: каждая работает на уровне \(I_v\) или одной ёмкости; ни одна не дотягивает до \(\eta_v\) как критерия замкнутой петли самооплаты.

\subsection{Assembly theory}\label{assembly-theory}

Marshall et al.~(2017) формализовали меру Pathway Complexity --- минимальное число рекурсивных шагов сборки объекта с переиспользованием частей; экспериментальное измерение через масс-спектрометрию сложных молекул и эмпирический порог (assembly index \(a\) порядка 15 для биогенных молекул) установлены в Marshall et al.~(2021), а развёрнутая программа сформулирована Sharma et al.~(2023). Объект с \(a\) выше этого порога интерпретируется как биосигнатура: глубина накопленной селекции, требуемая для сборки такого объекта, не достигается абиотическими процессами в наблюдаемых масштабах времени.

Assembly index работает на знаменателе диады «модель × ставка» --- мере накопленной селекции, --- но не определяет числитель. Assembly theory одинаково присваивает высокий \(a\) окаменелости и живой клетке, не различая накопленную в материи селекцию и актуальное удержание \(I_{\text{pred}}\). Коэффициент возврата ошибки в текущем поведении объекта в схеме не фигурирует. На шкале \(\eta_v\) assembly index занимает позицию операционального индекса сложности, не витальной шкалы: высокий \(a\) --- необходимое условие положительности \(\eta_v\) для систем со сложной молекулярной архитектурой, но не достаточное. Assembly index --- операциональный прокси для \(C_v\) или \(S_v\) в составе \(I_v\), не самостоятельная шкала.

\subsection{\texorpdfstring{Интегрированная информация \(\Phi\)}{Интегрированная информация \textbackslash Phi}}\label{ux438ux43dux442ux435ux433ux440ux438ux440ux43eux432ux430ux43dux43dux430ux44f-ux438ux43dux444ux43eux440ux43cux430ux446ux438ux44f-phi}

Тонони и его школа (Tononi 2004; Tononi et al.~2016) определили интегрированную информацию \(\Phi\) как меру неразложимости целого в сумму частей. В формализации IIT 3.0 (Oizumi et al.~2014) \(\Phi > 0\) интерпретируется как достаточное условие минимального феноменального сознания. В контексте демаркации витального \(\Phi\) --- мощный кандидат на формализацию информационной ёмкости \(C_v\): системы с высокой интегрированной информацией удовлетворяют требованию модели, в которой части не разложимы в независимые подсистемы.

Программа Тонони не вводит ограничения на термодинамическую плату удержания этой интегрированной информации и не требует, чтобы плата принадлежала самой системе. Бытовой термостат, по слабым формулировкам IIT, имеет ненулевое \(\Phi\) --- отсюда обсуждаемые в Tononi et al.~(2016) выводы о минимальном «опыте» простых систем. \(\Phi\) --- необходимое условие положительности \(C_v\), не достаточное условие положительности \(\eta_v\): термодинамическая принадлежность платы системе остаётся вне аппарата IIT. Связь продуктивна как декомпозиция: \(\Phi\) операционализирует \(C_v\), \(\eta_v\) добавляет требование \(T_v > 0\) с принадлежностью к той же системе.

\subsection{Телеодинамика}\label{ux442ux435ux43bux435ux43eux434ux438ux43dux430ux43cux438ux43aux430}

Дикон в \emph{Incomplete Nature} (Deacon 2011) построил иерархию homeodynamics → morphodynamics → teleodynamics. Третья ступень определяется как сцепление двух морфодинамических систем, взаимно ограничивающих саморазрушение друг друга; введено понятие \emph{absential causality} --- причинения тем, чего ещё нет (виртуальной целью, не данной актуально). Телеодинамика --- наиболее близкая к диаде «модель × ставка» программа: модель и ставка присутствуют в схеме логически, обе оси проговорены.

Дефект программы --- отсутствие операциональной количественной меры. Иерархия Дикона формулирована в качественных философских терминах, и реконструкция \(\eta_v\) или коэффициента возврата ошибки через absential causality технически невозможна без существенной формализации, которой текст программы не предлагает. Телеодинамика --- структурный портрет, в котором обе оси диады присутствуют, но количественная мера, проверяемая на бактерии или мегаполисе, отсутствует. Телеодинамика --- концептуальный предшественник, операционализация которого через шкалу \(\eta_v\) и есть один из вкладов настоящей работы.

\subsection{Принцип свободной энергии}\label{ux43fux440ux438ux43dux446ux438ux43f-ux441ux432ux43eux431ux43eux434ux43dux43eux439-ux44dux43dux435ux440ux433ux438ux438}

Принцип свободной энергии Фристона (Friston 2010) --- наиболее близкая к настоящей рамке формализация. В FEP-литературе различаются две величины. \emph{Variational free energy} (VFE)

\begin{equation}F[q, o] = D_{KL}\bigl[q(s) \,\|\, p(s|o)\bigr] - \ln p(o) \quad \bigl(= \mathbb{E}_q[\ln q(s) - \ln p(o, s)]\bigr), \tag{3}\end{equation}

или эквивалентно \(F = D_{KL}[q(s) \| p(s)] - \mathbb{E}_q[\ln p(o|s)]\) (complexity minus accuracy, \(-\text{ELBO}\)). Минимизируется в перцепции: распределение \(q(s)\) над скрытыми состояниями среды приближает истинное апостериорное \(p(s|o)\); \(-F\) --- нижняя граница на log-evidence \(\ln p(o)\). \emph{Expected free energy} (EFE)

\begin{equation}G[\pi] = \mathbb{E}_{q(o', s'|\pi)}\bigl[\ln q(s'|\pi) - \ln p(o', s')\bigr], \tag{4}\end{equation}

минимизируется в выборе действия / active inference. Здесь \(p(o', s')\) --- распределение предпочтений (генеративная модель с предпочтениями над будущими исходами и состояниями), и \(G\) канонически раскладывается на эпистемическую ценность (информационный выигрыш о скрытых состояниях) и прагматическую ценность (близость к предпочтениям). Различие фундаментально: \(F\) описывает обновление убеждений, \(G\) --- выбор действий (Friston 2010, 2019). При гауссовой generative model в стационарном режиме \(-F \lesssim I_{\text{pred}}\) --- концептуальное соответствие, связывающее шкалу и FEP в перцептивной части; формальная связь --- открытая техническая задача.

Каноническая FEP-as-life линия онтологизирует Markov blanket как границу живого. Канонические работы --- Kirchhoff et al.~(2018) «The Markov blankets of life»; Constant et al.~(2018); Ramstead et al.~(2020). BDDB-критика (Bruineberg et al.~2022) нацелена именно на эту линию.

Принципиальное возражение к FEP сформулировано в Bruineberg et al.~(2022, далее BDDB) --- target article в BBS-формате с открытым peer commentary. \emph{Markov blanket} (Pearl 1988) --- статистическое условие экранирования --- допускает два прочтения: инструментальное (\emph{Pearl-blanket}) и реалистское --- онтологизацию как объективной границы живой системы (\emph{Friston-blanket}, Friston 2019). Центральный тезис BDDB: Pearl-blanket и Friston-blanket неотличимы без скрытого ввоза эпистемических допущений; переход от первого ко второму квалифицируется как подмена понятий. Регулятор Уатта у BDDB --- диагностический пример, демонстрирующий, что FEP-аппарат не несёт собственной онтологической нагрузки (термостат как тривиальный active inference --- Baltieri and Buckley 2019). Параллельная критика --- Aguilera et al.~(2022, формальный анализ примеров без стационарной плотности) и Bruineberg et al.~(2018, экологически-энактивная линия). Дискуссия о парируемости BDDB-критики остаётся открытой.

Наш ход --- введение требования self-payment --- не реализует переход от Pearl- к Friston-blanket через термодинамический канал. \(\eta_v\) предлагает \textbf{альтернативную онтологию}: вместо статистической реальности blanket-разделения --- термодинамическую принадлежность диссипации одной физической системе. Self-payment --- допущение, независимое от BDDB-критики: сдвигает вопрос с «существует ли Friston-blanket» на «оплачивается ли удержание модели собственной диссипацией». \(\eta_v\) не наследует эпистемическую структуру, критикуемую BDDB. Термодинамическую принадлежность несёт предикат \(S\): обязательную диссипацию удержания модели оплачивает та же физическая граница \(X\) (counterfactual-ablation \(-dF_X/\partial t\), § 2.2, § 2.7), а не статистическое разделение скрытых состояний. Допущения (кандидатная граница, порог \(\delta\), окно \(\tau_d\)) эксплицитны и проверяемы через чувствительность (§ S2.4). У термостата predictive information о собственной среде мала, но \textbf{положительна} (биметалл отслеживает автокоррелированную температуру комнаты, \(I_{\text{pred}} \approx 3\) бит, см. ниже); демаркацию несёт не малость числителя \(\eta_v\), а провал самооплаты (\(S = 0\)): обязательную диссипацию оплачивает электросеть и петля возврата ошибки разомкнута --- независимо от существования у термостата Friston-blanket.

\emph{Counterfactual ablation термостата.} Применение процедуры § 2.2 к биметаллическому термостату с гистерезисом \(\Delta T_{\text{hyst}} = 1\) К, контролирующему комнату с тепловой инерцией \(\tau_{\text{room}} \sim 30\) мин, перебирает пять кандидатов на включение в границу:

\begin{itemize}
\tightlist
\item
  биметаллическую пластину (датчик и исполнитель);
\item
  контактную группу реле (коммутация энергии);
\item
  нагревательный элемент;
\item
  электросеть как источник тока;
\item
  внешнего уставщика.
\end{itemize}

Counterfactual-вычитание оставляет внутри границы первые два --- их удаление разрывает контур регулирования. Три остальных опознаются как внешние: вычитание нагревателя устраняет тепловой источник, не петлю принятия решений; вычитание электросети аналогично; вычитание уставщика не нарушает регулирование вокруг фиксированной уставки. Внутри границы \(E_{\text{actual}}\) --- упругая релаксация металла плюс диссипация реле, \(\sim 10^{-6}\) Вт (воспроизводимый расчёт двух интерпретаций --- тривиальной и управляющей --- в Supplementary Materials). Внутренний канал, удерживающий \(I_{\text{pred}}\) о температуре комнаты, --- деформация биметалла. Оценка: \(I_{\text{pred}} \le \log_2(\Delta T_{\text{room}}/\Delta T_{\text{hyst}}) \approx 3\) бит при суточном \(\Delta T_{\text{room}} \sim 10\) К (до \(\sim 5\) бит при сезонном размахе \(\sim 30\) К). Эта \(I_{\text{pred}}\) относится к среде «комната», не к среде «термостат + комната + источник питания». Петля принятия решений, оплачиваемых собственным распадом, у термостата неполна: нагрев комнаты оплачивается внешней электросетью, не собственной свободной энергией биметалла. \(\eta_v^{\text{thermostat}}\) не определён в обеих самосогласованных интерпретациях, но через провал разных ёмкостей. В управляющей задаче «регулировать температуру комнаты» проваливается \(T_v^{\text{int}}\): нагрев оплачивает электросеть, \(E_{\text{actual}}^{\text{own}} = 0\) относительно этой задачи. В тривиальной задаче «удерживать собственную форму биметалла» проваливается \(S_v^{\text{int}}\): пассивное thermo-mechanical coupling не образует петли возврата ошибки --- при отклонении формы биметалла никакая собственная диссипация её не корректирует. Расхождение между интерпретациями указывает не на онтологический переход, а на конвенциональность выбора границы как методологического инструмента: одна и та же физическая система допускает несколько самосогласованных границ. Bruineberg et al.~(2022) формально показывают, что такого онтологического перехода без скрытого ввоза эпистемических допущений сделать нельзя. Полный покомпонентный разбор пяти кандидатов --- Supplementary § S4.4.

Соотношение классов FEP-применимых и \(\eta_v\)-витальных систем асимметрично. Утверждение «FEP --- частный случай \(\eta_v\)» неверно по направлению. Правильное соотношение: шкала \(\eta_v\) выделяет внутри класса FEP-применимых систем подкласс, в котором минимизируемая \(F\) оплачивается собственной диссипацией. Вне этого подкласса --- термостат, маятник с внешним питанием, любой саморегулирующийся контур с экстернализованной платой --- проваливается \textbf{самооплата} (\(S = 0\)): минимизируемая \(F\) оплачивается извне, даже когда числитель \(\eta_v\) положителен. FEP сохраняет описательную силу на всём классе, но теряет дискриминационную силу: она не отделяет системы со ставкой от систем без ставки. \(\eta_v\) предъявляет термодинамический критерий, операционально независимый от FEP-формализма. Системы, удовлетворяющие FEP с непустой петлёй действия--восприятия и нетривиальной собственной диссипацией, как правило проходят \(\eta_v\)-верификацию. Системы, удовлетворяющие FEP в чисто формальной записи с экстернализованной энергией (термостат), \(\eta_v\)-верификацию проваливают.

\emph{Active inference и пространство соответствий.} Active inference (Friston et al.~2017; Parr et al.~2022) --- конкретная вычислительная схема, выводимая из FEP при выборе политики через EFE. В терминах активного вывода числитель \(\eta_v\) связан с predictive information в смысле generative model, а знаменатель --- с \textbf{информационной ценой удержания} generative model в собственных степенях свободы (\(I_{\text{mem}}\)); физическая цена удержания входит отдельным термодинамическим якорем (бонд Still), не в знаменатель \(\eta_v\). Это посылка \(\eta_v\)-программы, добавляемая к стандартной AIF, --- не выводится из EFE-минимизации напрямую. Прагматическая и эпистемическая ценность политик связаны с числителем (predictive information) через структуру generative model, не со знаменателем напрямую. Развёрнутое сопоставление вокабуляров (action, policy, прагматическая ценность, эпистемическая ценность против \(I_{\text{pred}}\), \(E_{\text{actual}}\), \(T_v\), \(S_v\)) --- задача отдельной работы; здесь существенно, что соответствие технически возможно и не требует пересмотра ни одной из исходных рамок.

\emph{Andrews 2021 и Williams 2020.} Andrews (2021) атакует FEP за невыводимость эмпирического содержания из формализма; критика парируется в \(\eta_v\) привязкой числителя и знаменателя к измеримым величинам. Williams (2020) атакует predictive coding как нейронный механизм, не predictive information как формальную меру; критика нейтральна к \(\eta_v\). Развёрнутый разбор и различение predictive coding (Rao and Ballard 1999; Clark 2013) против predictive information (Bialek et al.~2001) --- Supplementary § S4.4a.

Сводное положение: формализм минимизации свободной энергии становится содержательным критерием витальности только при условии self-payment. Без этого условия FEP описывает любой саморегулирующийся контур и теряет дискриминационную силу. С этим условием FEP уточняется до своей \(\eta_v\)-дисциплины. Для систем, где FEP применим с непустой петлёй действия--восприятия и нетривиальной собственной диссипацией, \(\eta_v\) и FEP дают согласующиеся характеристики. Для систем, где FEP применим лишь в чисто формальной записи с экстернализованной энергией, проваливается самооплата (\(S = 0\)): демаркацию несёт предикат \(S\), а не обнуление \(\eta_v\).

\subsection{Демаркация от соседних программ: сводная диагностика}\label{ux434ux435ux43cux430ux440ux43aux430ux446ux438ux44f-ux43eux442-ux441ux43eux441ux435ux434ux43dux438ux445-ux43fux440ux43eux433ux440ux430ux43cux43c-ux441ux432ux43eux434ux43dux430ux44f-ux434ux438ux430ux433ux43dux43eux441ux442ux438ux43aux430}

Сводная позиционная диагностика основных программ сведена в таблице ниже и выявляет единый структурный паттерн дефекта: каждая программа берёт всерьёз ровно одну сторону витальности и оставляет остальные за пределом аппарата. Линия organisational closure (Maturana--Varela, Rosen, Mossio--Montévil--Longo) и линия программ OoL развёрнуто разобраны в § 4.7 и § 4.8 соответственно; в таблице они присутствуют как однострочные позиционные сводки.

\begin{center}\small
\begin{tabular}{@{} >{\raggedright\arraybackslash}p{0.1594\linewidth} >{\raggedright\arraybackslash}p{0.3043\linewidth} >{\raggedright\arraybackslash}p{0.5362\linewidth}@{}}
\toprule
\begin{minipage}[b]{\linewidth}\raggedright
Программа
\end{minipage} & \begin{minipage}[b]{\linewidth}\raggedright
Ось, взятая всерьёз
\end{minipage} & \begin{minipage}[b]{\linewidth}\raggedright
Что оставлено за пределом аппарата
\end{minipage} \\
\midrule
Walker--Cronin (assembly theory) & знаменатель диады --- накопленная селекция в материи (assembly index \(a\)) & числитель имплицитен; \(a\) одинаков у окаменелости и живой клетки --- витальная дискриминация требует внешней интуиции о петле возврата ошибки \\
Tononi (IIT) & мера информационной интеграции \(\Phi\) & термодинамическая плата; \(\Phi\)-положительный термостат --- известная критическая точка теории \\
Deacon (телеодинамика) & структурная диада (absential causality) & количественная мера: absential causality не реконструируется через измеримые величины \\
Friston (FEP) & формализм минимизации свободной энергии & требование принадлежности модели и диссипации одной системе; устойчивая критика Bruineberg et al.~2022 ставит под вопрос путь онтологизации Friston-blanket (открытая дискуссия без устоявшегося консенсуса), \(\eta_v\) предлагает альтернативную термодинамическую онтологию, не зависящую от исхода этой дискуссии (см. § 4.4) \\
Maturana--Varela / Rosen / Mossio--Montévil--Longo (organisational closure) & требование принадлежности петли одной системе & количественная термодинамическая дисциплина: качественное закрытие без термодинамической дисциплины бонда (2) (см. § 4.7) \\
Hordijk--Steel / England / Maynard Smith--Szathmáry / Pross / Smith--Morowitz (OoL программы) & операционализируют одну сторону abiotic→biotic перехода каждая & объединяющая операциональная мера self-payment как единый термодинамический критерий момента origin of life (см. § 4.8) \\
\bottomrule
\end{tabular}
\end{center}

Перечисленные дефекты не случайны: они отражают разные способы взять одну сторону витальности всерьёз и оставить остальные за пределом аппарата. \(\eta_v\) как отношение, требующее принадлежности обоих компонентов одной системе, --- операционализация, не позволяющая ни одной из четырёх осей раствориться в фоне.

\subsection{Связь с традицией: краткая позиционная сводка}\label{ux441ux432ux44fux437ux44c-ux441-ux442ux440ux430ux434ux438ux446ux438ux435ux439-ux43aux440ux430ux442ux43aux430ux44f-ux43fux43eux437ux438ux446ux438ux43eux43dux43dux430ux44f-ux441ux432ux43eux434ux43aux430}

\(\eta_v\)-программа располагается на пересечении классических традиций. Линия Schrödinger (1944, aperiodic crystal + negative-entropy-formula) и Szilard--Landauer--Bennett (1929--1989, термодинамическая стоимость информационных операций) --- непосредственное наследование числителя--знаменателя. Неравновесная термодинамика Пригожина--Николиса (Nicolis and Prigogine 1977) --- предшественник \(T_v\). Wiener--Ashby--Foerster (Wiener 1948; Ashby 1952; von Foerster 2003) --- кибернетика второго порядка. Сюда же примыкают anticipatory systems Розена (Rosen 1985) и skin-in-the-game Талеба (Taleb 2018, концепт ставки, формализованный в § 2.1). Алгоритмическая информация Колмогорова (Kolmogorov 1965) --- предшественник \(C_v\). Автокаталитические замыкания Кауфмана (Kauffman 1993, 2000) и программная рамка Walker--Davies--Ellis (Walker et al.~2017a) дополняют список. \(\eta_v\) --- не наследник какой-либо одной из них, а юкстапозиция: термодинамическая рамка плюс предсказательное содержание плюс структурный счёт ёмкостей через counterfactual ablation (§ 2.2). Шкала собирает обе оси диады модели/ставки в одну измеримую величину и явно вводит требование собственной диссипации. Каждая классическая программа операционализируется через одну из ёмкостей в составе \(I_v\); объединение в одну витальную шкалу --- самостоятельный шаг, который ни одна из исходных рамок не делала. Расширенная позиционная карта --- § S4.7--S4.9 supplementary.

\subsection{Линия организационного замыкания}\label{ux43bux438ux43dux438ux44f-ux43eux440ux433ux430ux43dux438ux437ux430ux446ux438ux43eux43dux43dux43eux433ux43e-ux437ux430ux43cux44bux43aux430ux43dux438ux44f}

\emph{Линия организационного замыкания} (autopoiesis Maturana--Varela 1980, M,R-systems Rosen 1991, constraint closure Mossio--Montévil--Longo 2016, organisational closure Bich--Green 2018, плюс epistemic-cut Pattee 1982) формулирует жизнь через замкнутость каузальных процессов. Замыкание берётся относительно efficient causation или constraint network, производящих компоненты системы. \(\eta_v\)-программа категориально отлична: organisational closure формализует производство компонентов и качественную циркулярность; self-payment формализует термодинамическую принадлежность бюджета модели одной физической системе с моделирующим. \emph{Дифференциальный тест} на паре \emph{E. coli} / JCVI-syn3.0 (Hutchison et al.~2016) различает шкалы внутри одного класса организационного замыкания: качественные критерии не различают эти системы, \(\eta_v\) ранжирует их через произведение информационной и термодинамической компонент. Self-payment не претендует на ребрендинг или замену organisational closure --- он добавляет к её качественной концептуальной работе количественную термодинамическую дисциплину через бонд Still (2). Полный разбор четырёх линий (Pattee, Maturana--Varela, Rosen, Mossio--Montévil--Longo) с категориальными различениями и дифференциальным тестом --- § S4.7 supplementary.

\subsection{Линия происхождения жизни}\label{ux43bux438ux43dux438ux44f-ux43fux440ux43eux438ux441ux445ux43eux436ux434ux435ux43dux438ux44f-ux436ux438ux437ux43dux438}

\emph{Линия происхождения жизни} представлена в Theory in Biosciences и смежных журналах пятью относительно автономными традициями. Гиперцикл Эйгена--Шустера (Eigen 1971; Eigen and Schuster 1979) --- replicator-level автокаталитическое замыкание с template-носителем. RAF-сети (Hordijk and Steel 2004) --- формализм коллективной автокаталитичности. Dissipative adaptation (Perunov, Marsland, and England 2016; статистическая основа --- England 2013, \emph{Statistical physics of self-replication}) --- термодинамический предшественник якоря самооплаты (бонд Still) и ёмкости \(T_v\), не знаменателя \(\eta_v\). Major transitions (Maynard Smith and Szathmáry 1995; Szathmáry 2015) --- иерархия эволюционных уровней с переадресацией субъекта. Геохимические программы (Pross 2012; Smith and Morowitz 2016; Kauffman 2019; Martin and Russell 2003) --- обнаружение замкнутой петли поверх геохимического потока. Каждая операционализирует одну сторону перехода abiotic→biotic. В \(\eta_v\)-терминах все пять артикулируют один и тот же фазовый переход. Момент origin of life --- момент, когда геохимическая система переходит от внешней оплаты к self-payment: часть свободной энергии замыкается внутри собственного контура и физически оплачивает удержание модели среды. Положительность \(\eta_v\) --- \textbf{необходимое условие} этого момента в каждой из программ; достаточные условия включают дарвиновскую динамику на популяционном уровне (§ 2.7). Операциональная процедура измерения \(\eta_v\) на pre-biotic геохимическом субстрате не определена и составляет открытую задачу программы. Развёрнутый разбор пяти традиций с указанием точек контакта и параллельной программы Chirumbolo--Vella (2026) --- § S4.8 supplementary.

\subsection{Позиционирование в современных демаркационных дебатах}\label{ux43fux43eux437ux438ux446ux438ux43eux43dux438ux440ux43eux432ux430ux43dux438ux435-ux432-ux441ux43eux432ux440ux435ux43cux435ux43dux43dux44bux445-ux434ux435ux43cux430ux440ux43aux430ux446ux438ux43eux43dux43dux44bux445-ux434ux435ux431ux430ux442ux430ux445}

\(\eta_v\)-программа методологически совместима с шестью устойчивыми линиями современной философии биологии. Это anti-definitionism Cleland (2019), life-as-lineage Mariscal--Doolittle (2020), operational-frontiers Bich--Green (2018), process-biology Pradeu (2018), информационно-теоретическая индивидуальность Krakauer et al.~(2020) и basic-autonomy Ruiz-Mirazo--Moreno (2004). Программа выступает не как альтернатива им, а как формальная количественная шкала, на которой эти линии могут совместно работать. \(\eta_v\) не предлагает очередное определение жизни через списки признаков, но вносит операциональный дискриминатор с явным фальсификационным режимом (§ 5) и процессуальной онтологией (шкала вычисляется на окне \(\tau_d\), § 2.6). На пограничных случаях (вирус-в-изоляции, минимальная клетка, LLM-как-корпорация) шкала ведёт себя в согласии с биологической конвенцией через counterfactual ablation § 2.2. Это вклад в продолжающуюся демаркационную дискуссию, не предложение нового определения «что есть жизнь». Развёрнутая позиционная диагностика по каждой из шести линий --- § S4.9 supplementary.

\section{Критерии фальсификации}\label{ux43aux440ux438ux442ux435ux440ux438ux438-ux444ux430ux43bux44cux441ux438ux444ux438ux43aux430ux446ux438ux438}

Программа перечисляет ниже девять проверяемых следствий в трёх блоках. Реальная фальсификационная поверхность, однако, уже (§ 2.7, «Асимметрия фальсификации»): ложно-положительные случаи исключены конструкцией, и эмпирический риск концентрируется в ложно-отрицательном проникновении (Блок 2 п. 3) и провале прогрессивных предсказаний (Блок 3 п. 1, несущий Прогноз 1). Процедурное правило, общее для всех: граница системы и прокси \(I_{\text{pred}}\), \(I_{\text{mem}}\) фиксируются до верификации \(\eta_v\) и не пересматриваются ради получения желаемого знака (§ 2.2). Фальсификация работает на двух уровнях --- методологическом (дефект текущей операционализации) и онтологическом (несостоятельность самого концепта); различение фиксируется post hoc для каждой реализовавшейся линии.

\subsection{Блок 1 --- Прямые опровержения шкалы}\label{ux431ux43bux43eux43a-1-ux43fux440ux44fux43cux44bux435-ux43eux43fux440ux43eux432ux435ux440ux436ux435ux43dux438ux44f-ux448ux43aux430ux43bux44b}

Атакуют структуру \(\eta_v\) через обнуление одной из трёх ёмкостей.

\begin{enumerate}
\def\labelenumi{\arabic{enumi}.}
\tightlist
\item
  \textbf{\(T_v\) провал (самооплата термодинамического пола).} Наблюдение: витально-зачтённая (\(S = 1\)) система, чья обязательная по Still диссипация удержания демонстрируемо поставляется \textbf{из-за} её counterfactual-ablation-границы (§ 2.2) --- пол (2) оплачивается извне. Опровергает: тезис самооплаты термодинамического пола (не сам бонд Still, который есть теорема стохастической термодинамики); обнуляет \(T_v\) как необходимое условие.
\item
  \textbf{\(C_v\) провал.} Наблюдение: витально-зачтённая система, удерживающая predictive information о дрейфующей среде на окне \(\tau_d\), у которой обязательная по бонду (2) цена удержания непредсказательной памяти \textbf{не приписывается} её собственному бюджету (при вердикте \(S = 1\) оплачивается извне). Опровергает: согласованность предиката самооплаты \(S\) с \(C_v\) как необходимым условием.
\item
  \textbf{\(S_v\) провал.} Наблюдение: система с \(\eta_v > 0\) при одновременных \(Q < Q_{\text{null}} + 2\sigma\) (модулярность не отличима от degree-preserving нуль-модели --- configuration model, сохраняющего распределение степеней) и \(C(k) \not\sim k^{-\beta_h}\) с \(\beta_h \in [0{,}5; 1{,}5]\) (отсутствие иерархической сигнатуры по Ravasz--Barabási 2003). Опровергает: \(S_v\) как необходимое условие.
\end{enumerate}

Опровержение любой ёмкости одновременно обрушивает обе шкалы --- вердикт \(V(X)\) (через \(\eta_v\), если провалена \(C_v\); через предикат \(S\), если \(T_v\) или \(S_v\)) и скрининговый индекс \(I_v\) через геометрическое среднее. Это одна фальсификация через ёмкость, не двойная.

\subsection{Блок 2 --- Тесты прокси и операционализации}\label{ux431ux43bux43eux43a-2-ux442ux435ux441ux442ux44b-ux43fux440ux43eux43aux441ux438-ux438-ux43eux43fux435ux440ux430ux446ux438ux43eux43dux430ux43bux438ux437ux430ux446ux438ux438}

Атакуют не структуру шкалы, а её привязку к измеримым величинам.

\begin{enumerate}
\def\labelenumi{\arabic{enumi}.}
\tightlist
\item
  \textbf{Нарушение монотонности темпа адаптации по \(\eta_v\).} Наблюдение: в контролируемой экспериментальной эволюции системы с \(\eta_v\), различающимся на порядок, дают темпы адаптации, неразличимые статистически на длинных окнах. Опровергает: центральное отношение лишается числового смысла (конструктивная форма --- Прогноз 1 ниже).
\item
  \textbf{Витальность по трём независимым каналам при \(\eta_v = 0\).} Наблюдение: система с положительной витальностью по темпу адаптации, скорости самовосстановления и биохимическим маркерам самопроизводства (Joyce 1994), но \(\eta_v = 0\) при всех санкционированных операционализациях. Опровергает: числитель как операциональный прокси (методологический уровень, § 2.7) либо концепт витальности через эффективность самомоделирования (онтологический).
\item
  \textbf{Конвенционально живая система с \(\eta_v = 0\).} Наблюдение: биологическая система, заведомо живая по конвенциональным критериям, со строго нулевым \(\eta_v\) при всех разумных операционализациях. Опровергает: адекватность шкалы как операционального дискриминатора витальности.
\end{enumerate}

\subsection{Блок 3 --- Прогрессивные предсказания}\label{ux431ux43bux43eux43a-3-ux43fux440ux43eux433ux440ux435ux441ux441ux438ux432ux43dux44bux435-ux43fux440ux435ux434ux441ux43aux430ux437ux430ux43dux438ux44f}

Избыточные эмпирические следствия программы, не выводимые из конкурирующих рамок.

\begin{enumerate}
\def\labelenumi{\arabic{enumi}.}
\tightlist
\item
  \textbf{Прогноз 1 --- степенная зависимость темпа адаптации.} Здесь \(\eta_v \equiv \eta_v^{\mathrm{do}}\) --- измеренная \textbf{для каждой линии} пассивная (интервенционная) свежесть модели: иммобилизованный FRET каждого штамма (\(\mathrm{do}(a) =\) иммобилизация, § 2.1), а \textbf{не} активная свежесть свободно плывущих клеток. Поэтому прогноз \textbf{корректно поставлен} (предиктор \(\eta_v^{\mathrm{do}}\) определён и измерим иммобилизованным FRET) независимо от открытого статуса активной границы \(\eta_v\): штаммы с большей пассивной свежестью модели адаптируются быстрее. При этом \textbf{принимается, что \(\eta_v^{\mathrm{do}}\) --- монотонный прокси активной предсказательной свежести}, гонящей адаптацию свободно плывущих клеток, и \textbf{адекватность этого прокси --- часть того, что Прогноз 1 проверяет/фальсифицирует}. Дополнительно принимается \textbf{межшкальный мост}: синхронная свежесть сигнального контура (секунды--минуты, внутри одной клетки) монотонно трекает скорость, с какой линия переподгоняет модель под отбором за поколения; этот переход внутриклеточный темп → межпоколенный --- тоже часть фальсифицируемого содержания Прогноза 1. Наблюдение: на 45 линиях \emph{E. coli} в хемостате \(r_{\text{adapt}} = A \cdot \eta_v^{\alpha}\) с \(\alpha \in [0{,}3;\, 1{,}0]\) и коэффициентом корреляции \(\rho \in [0{,}4;\, 0{,}7]\) (Пирсон между \(\log r_{\text{adapt}}\) и \(\log \eta_v\)). Опровергает: \(\hat{\alpha}\) значимо вне \([0{,}3;\, 1{,}0]\) при \(p < 0{,}01\) (неправильный показатель степени); \(\hat{\rho} < 0{,}2\) (нет detectable trend); \(\hat{\rho} < 0\) (тренд в противоположном направлении --- качественное опровержение).
\item
  \textbf{Прогноз 2 --- монотонность \(\eta_v\) по окну согласования.} Наблюдение: в режиме приближения к стационарности при нормированной диссипации ожидается монотонный рост \(\eta_v(\Delta\tau)\) на интервале \(\Delta\tau \ll \tau_d\). Опровергает: убывание \(\eta_v\) с \(\Delta\tau\) --- контр-аргумент к стационарной рамке как таковой.
\item
  \textbf{Прогноз 3 --- эмпирический порог интернализации энергобюджета \(f^*\)} (где \(f = E_{\text{actual}}^{\text{own}}/E_{\text{actual}}^{\text{total}}\) --- доля собственной энергии в полном бюджете). Структурный предикат замкнутости петли самооплаты \textbf{бинарен} (counterfactual ablation § 2.2: канал возврата ошибки либо самооплачивается, либо нет); \(f\) --- непрерывный прокси этого замыкания, не сама решающая граница. Наблюдение: для гипотетических автономных воплощённых агентов доля прошедших скрининг (\(\eta_v^{\text{int}} > 0\)) резко возрастает выше эмпирического порога \(f^* \in [0{,}1;\, 0{,}5]\) --- статистическая регулярность совместной интернализации (персистентное состояние + собственная оплата + физический канал возврата ошибки), а не структурная дефиниция. Опровергает: систематическое наблюдение системы с \(f \in [0; 0{,}1]\), проходящей скрининг, либо с \(f > 0{,}5\), проваливающей. Differential prediction относительно AI-safety (Hubinger et al.~2019; Carlsmith 2022): порог \emph{способностей} независим от \emph{термодинамического бюджета} \(f^*\); scope не applicable к нынешним cloud-hosted LLM без сохранения состояния (\(f = 0\) конструктивно).
\end{enumerate}

\emph{Условие подтверждения уникальности операционализации (не критерий фальсификации).} Обнаружение альтернативного прокси, дающего ту же демаркационную дискриминацию, но не сводимого к \(\eta_v\) через тождество или монотонное преобразование, не опровергает программу --- оно подтверждает содержательность демаркации.

Полные операциональные протоколы (численные критерии, p-values, выборки, требования предрегистрации, Bennett-режим caveat для Блока 1 п. 1, NASA-референт для Блока 2 п. 2, scope-сценарии Прогноза 3) --- § S5 supplementary. Сводный лакатошевский паспорт программы --- § S5.4.

\section{Обсуждение и открытые вопросы}\label{ux43eux431ux441ux443ux436ux434ux435ux43dux438ux435-ux438-ux43eux442ux43aux440ux44bux442ux44bux435-ux432ux43eux43fux440ux43eux441ux44b}

\emph{Астробиология.} Главное следствие --- переформулировка программы поиска жизни. Большинство объектов проходит \(I_v\)-скрининг (структурный потенциал универсально высок) и проваливает \(\eta_v\)-верификацию (актуальная замкнутость петли самооплаты редка). Операционально: при разработке миссий поиска биосигнатур (Europa Clipper, JUICE, Dragonfly, экзопланетная спектроскопия) приоритет смещается с химической неравновесности на поиск тройной одновременности --- неравновесной химии (\(T_v\)), биогенно-подобной информационной подписи (\(C_v\), через изотопное фракционирование и хиральность), иерархически-модулярной топологии среды (\(S_v\)). Каждое --- необходимое условие; одновременность всех трёх --- операциональный кандидат на биосигнатуру. Связь с уже существующей методологической рамкой Ladder of Life Detection (Neveu et al.~2018): \(\eta_v\) позиционируется как один из срезов operational criteria для life detection, дополняющий химико-структурные ступени лестницы термодинамическим требованием self-payment.

\emph{Происхождение жизни как замыкание петли самооплаты.} Программа поиска биосигнатур фиксирует \textbf{где} жизнь существует; параллельный вопрос --- \textbf{когда} она возникает на эволюционирующем геохимическом субстрате. Здесь, как и в § 3.4 для AI-систем, требуется явное двухуровневое различение.

\emph{Структурный уровень.} Положительность \(\eta_v\) --- \textbf{необходимое условие} момента origin of life (согласовано с § 4.8 и формулировкой «необходимого, не достаточного» термодинамического требования self-payment). Достаточные условия включают дарвиновскую динамику (§ 2.7), вошедшую в работающее NASA-определение жизни как «self-sustaining chemical system capable of Darwinian evolution» (Joyce 1994); \(\eta_v > 0\) фиксирует термодинамический срез, эволюционная динамика --- селективный. Структурно момент прохождения counterfactual ablation § 2.2 --- момент, когда геохимическая сеть впервые конституируется как замкнутая каталитическая петля с self-payment: часть свободной энергии перестаёт проходить транзитом и начинает оплачивать удержание собственной predictive information о среде. До этого момента \(\eta_v\) \textbf{не определён} (нет петли как объекта измерения); после --- измерим непрерывно. Граница «до» и «после» --- бинарный структурный предикат counterfactual ablation, не градиент внутри положительного класса.

\emph{Эмпирический уровень.} Операциональная процедура измерения \(\eta_v\) на pre-biotic геохимическом субстрате не определена и составляет открытую задачу программы --- программное обещание / направление исследований, требующее разработки proxies для \(I_{\text{pred}}\) и self-payment в эволюционирующей геохимической сети. Counterfactual ablation границы § 2.2 и MI-оценка \(I_{\text{pred}}\) через сенсорный канал применимы к уже-сформированным биологическим системам; их перенос в pre-biotic регистр --- конкретное направление операционального расширения. Возможные шаги: (а) RAF-сети (Hordijk and Steel 2004; Hordijk 2019) как структурный proxy замкнутости каталитической петли --- counterfactual subtraction отдельных компонентов сети есть операциональная реализация § 2.2 на геохимическом графе. (б) Изотопные сигнатуры (фракционирование \(\delta^{13}\mathrm{C}\), \(\delta^{34}\mathrm{S}\)) и хиральные подписи --- proxy для \(C_v\) (память сети о состоянии среды). (в) интегрированный термодинамико-сетевой тест на положительность self-payment через измеренный свободно-энергетический баланс RAF-замкнутой подсети относительно её reactive окружения. Разработка таких proxies и их калибровка на современных биоминиатюрах (минимальные клетки, синтетические протоклетки) --- будущая работа.

Связь с гидротермальными гипотезами --- прямая. Сюда относятся alkaline-vent сценарии Martin and Russell (2003), метаболические core программы Smith and Morowitz (2016), dynamic kinetic stability Pross (2012), adjacent-possible Кауфмана (Kauffman 2019); общая физическая рамка OoL --- Walker (2017b). Каждая из этих программ артикулирует тот же структурный переход со своей стороны, без объединяющей термодинамической меры. Конструктивная диагностика формулируется как направление разработки, а не как уже доступный протокол. Она требует одновременного появления трёх компонентов: каталитической замкнутости через RAF-анализ геохимической сети (Hordijk and Steel 2004; Hordijk 2019), информационной памяти через изотопные и хиральные подписи, положительного термодинамического счёта по схеме § 5. Это операциональная реформулировка задачи астробиологии: не только «найти, где жизнь есть», но и «детектировать структурный переход к ней» --- фальсифицируемая в случае обнаружения геохимических систем с устойчивым \(\eta_v > 0\) без характерной триады, как только методология будет операционально доразвёрнута.

\emph{Этика искусственного интеллекта.} Переход от вопроса «достаточно ли модель сложна, чтобы быть моральным субъектом» к вопросу «начнёт ли модель нести собственные ошибки сама». Появление собственной диссипативной платы за модель, удерживаемую в собственных степенях свободы работающего контура, есть структурное условие, при котором \(\eta_v\) выходит из области неопределённости в положительные значения. Фазовость относится к определению петли самооплаты (структурный предикат counterfactual ablation § 2.2 --- бинарен), не к эмпирическому переходу системы; эмпирическая динамика появления петли в реальной популяции AI-систем --- открытый вопрос с ожидаемо непрерывным профилем по доле интернализованной оплаты. Текущая стандартная архитектура LLM (прямой проход без сохранения состояния с экстернализованной оплатой инфраструктуры) исключает выполнение структурного условия. Путь к положительному \(\eta_v^{\text{int}}\) для искусственной системы требует одновременной интернализации (а) персистентное состояние, (б) собственной оплаты вычислений, (в) физического канала возврата ошибки --- не реализованных совместно ни в одной существующей системе на момент написания. Речь идёт о структурном условии, не о календарном прогнозе. Формальная привязка к триаде. (а) Персистентное состояние операционализирует \(C_v^{\text{int}} > 0\) (канал памяти для удержания). (б) Собственная оплата вычислений --- \(T_v^{\text{int}} > 0\) (собственный термодинамический бюджет). (в) Физический канал возврата ошибки --- \(S_v^{\text{int}} > 0\) (замкнутость топологии петли). Одновременность всех трёх даёт \(\eta_v^{\text{int}} > 0\).

«Нести собственные ошибки» --- в техническом смысле § 2.1: ошибки модели возвращаются физической потерей свободной энергии самой системы, не репутационными санкциями через внешний корпоративный субстрат. Падение частоты использования LLM и её переобучение силами корпорации --- плата корпорации, входящая в \(\eta_v\) корпорации; \(\eta_v^{\text{int}}\) модели как агента остаётся нулевым. Штрафы регулятора, юридические санкции против разработчика, рыночные репутационные потери --- возвращают ущерб не системе, удерживающей модель, а её внешнему субстрату.

\emph{Сложные системы и социальная динамика.} В применении к городам и цивилизациям шкала открывает два порога: структурный скрининг (\(I_v\)) и верификация (\(\eta_v\)). Цивилизация с высоким \(I_v\) и проседающим \(\eta_v\) --- структурно богатая, проваливающая удержание петли --- диагностически отличается от цивилизации с проседающим скринингом (разрушение инфраструктуры). Дифференциация полезна для исследований устойчивости социальных систем; требует эмпирической верификации.

\emph{Парадокс Ферми и фильтр экстернализации.} Гипотеза paper'а в применении к парадоксу Ферми: молчание Вселенной объясняется не редкостью жизни, а \textbf{схлопыванием технологических цивилизаций на пороге экстернализации энергобюджета}. В терминах рамки: цивилизация проходит «фильтр» тогда и только тогда, когда её \(\eta_v\) положителен через собственную диссипацию (\(T_v\) удержан внутри её контура). Технологическое развитие, направленное на использование внешних источников (электросеть → реакторы → мегаструктуры), увеличивает \(E_{\text{actual}}^{\text{external}}\) без увеличения \(E_{\text{actual}}^{\text{own}}\) --- большинство цивилизаций проходит \(I_v\)-скрининг (видимая структурная сложность: радиосигналы, мегаструктуры), но \textbf{проваливает \(\eta_v\)-верификацию} по тому же механизму, что LLM без сохранения состояния при границе «веса в момент инференса». Такие цивилизации нестабильны на эволюционных временах и не сохраняются как стационарные SETI-источники. Утверждение имеет статус спекулятивной экстраполяции и подлежит независимой эмпирической проверке. Проверяемое статистическое предсказание о наблюдаемых SETI-сигналах: транзиентные, неповторяющиеся, степенное распределение по длительности с обрывом --- индикатор фильтра; стационарный спектр с ИК-избытком, удерживающийся на десятилетних окнах, --- индикатор прохождения обоих порогов. Существующие null-результаты Dyson-sphere поисков (G-HAT survey по \textasciitilde{}\(10^5\) WISE-галактик, Griffith et al.~2015; обоснование программы --- Wright et al.~2014) согласуются с предсказанием фильтра, но не опровергают альтернативные сценарии (rare-Earth-class hypotheses, технозрелость ниже астрофизического разрешения). Условие опровержения: обнаружение трёх и более независимых стационарных SETI-источников с подтверждённой ИК-сигнатурой мегаструктур на интервалах \(\gtrsim 50\) лет несовместимо с предсказанием фильтра экстернализации в его текущей формулировке.

\emph{Статус следствия.} Переформулировка имеет статус спекулятивной экстраполяции: ни NASA Astrobiology Strategy, ни ESA Cosmic Vision, ни планетарный Decadal Survey не оперируют категорией «петли самооплаты». Если стандартные программы достигнут надёжной демаркации внеземной жизни без обращения к \(\eta_v\), рамка не принесёт астробиологически новой различительной силы и редуцируется к существующим агностическим биосигнатурам.

\emph{Открытые вопросы.} Нестационарное обобщение стационарного результата о сносе памяти (§ 2.1) для систем с растущей памятью --- биосферы на эоновых окнах, обучающихся искусственных систем --- требует строгого доказательства с учётом стоимости расширения памяти (разрабатывается в сопровождающей работе, Andriishin 2026). Каталог операциональных прокси для \(I_{\text{pred}}\) требует расширения с числовыми примерами и границами применимости. Случай Геи --- вопрос о единстве модели биосферы --- открыт эмпирически и зависит от данных о согласованности биогеохимических циклов на эоновых окнах. Операциональный протокол для обнаружения витального ИИ требует прецизионных порогов для популяционных экспериментов. Философский статус \(\eta_v\) --- позиция умеренного операционализма с разграничением онтологического референта и операциональной меры --- задача отдельной работы.

\emph{Dormant states.} Биологические состояния обратимой остановки --- споры, тихоходка в крио-анабиозе, замороженные эмбрионы, лизогенный цикл фага --- формально дают \(\dot{I}_{\text{pred}} = 0\) и \(E_{\text{actual}} \to 0\), что выводит \(\eta_v\) из области определения шкалы (§ 2.7). Согласование с биологической конвенцией («споры живы») достигается через окно усреднения \(\tau_d\): при выборе \(\tau_d\) на эволюционном масштабе (полный цикл sporulation \(\to\) активность \(\to\) sporulation) знак \(\eta_v\) возвращается к положительному. Формальное расширение шкалы на режимы временной обратимости знака --- направление сопровождающей работы (Andriishin 2026).

\emph{Граничные условия применимости.} Шкала \(\eta_v\) не делает трёх типов утверждений.

\begin{enumerate}
\def\labelenumi{\arabic{enumi}.}
\tightlist
\item
  Не утверждает тождество витальности и сознания. Высокий \(\eta_v\) --- необходимое условие сложной когнитивной деятельности (без удержания собственной predictive information невозможны ни предсказание, ни планирование, ни рефлексия), но не достаточное. Вопрос о феноменальном опыте, обсуждаемый в рамках интегрированной информации (Tononi et al.~2016) и трудной проблемы сознания, лежит за пределами шкалы \(\eta_v\).
\item
  Не предполагает биоцентризма по обратному жесту. Высокое \(\eta_v\) возможно у систем небиологического субстрата --- крупных городов, биосфер планетарного масштаба, гипотетических искусственных воплощённых агентов с замкнутой петлёй самооплаты. Это содержательное предсказание, а не калибровочная подгонка.
\item
  Не сводится к моральному критерию. Цивилизация, проваливающая \(\eta_v\)-верификацию при сохранном \(I_v\)-скрининге, диагностируется как структурно богатая, но невитальная --- диагноз структурного состояния, не этический вердикт. Этические следствия программы --- отдельная тема, опирающаяся на структурный диагноз, но не выводимая из него прямой импликацией.
\end{enumerate}

\emph{Стационарный режим как явное ограничение области применимости.} Все приведённые результаты --- определения, демаркационные тесты на камертонах, условия фальсификации --- формулируются в предположении стационарности среды и конечности памяти системы. Это не упрощение по умолчанию: стационарность стабилизирует \(I_{\text{mem}}\) и делает пошаговую цену удержания (бонд Still (2)) стационарной, и потому позволяет резко формулировать демаркационные ответы. В нестационарном режиме --- с растущей памятью, дрейфующей средой и непустым архивным слоем --- формализм требует расширения: кумулятивной диссипативной цены, учёта стоимости удержания исторических битов и феномена информационной ностальгии. Это расширение --- содержание сопровождающей работы (Andriishin 2026). Там доказывается условная лемма о неизбежности ностальгии в каноническом классе медленного дрейфа (OU-параметризация с FIFO-refresh и \(c < 1\)). Количественная адиабатическая асимптотика \(\dot{\eta_v}^{\text{excess}}\) для систем без периодического сброса памяти формулируется там как открытая гипотеза, поддержанная параметрическим сканом. Информационная ёмкость \(C_v\) настоящей работы там методологически уточнена через разбиение \(C_v^{\text{static}}\) (statistical complexity \(C_\mu\), структурная память --- § 2.3) против \(C_v^{\text{predictive}}\) (excess entropy \(\mathbf{E}\), предсказательная часть --- § 2.3), связанных соотношением \(C_v^{\text{predictive}} = (1 - \nu) \cdot C_v^{\text{static}}\), где \(\nu\) --- доля ностальгических битов; в стационарном режиме настоящей работы расхождение \(C_v^{\text{static}}\) и \(C_v^{\text{predictive}}\) перестаёт расти (\(\dot\nu \to 0\), \(\nu\) выходит на постоянный уровень) и служит диагностикой дрейфа лишь в нестационарном расширении, что согласует обе работы. Применимость стационарного \(\eta_v\), развёрнутого в настоящей статье, ограничена системами с эргодической средой и конечной памятью; для биосфер на эоновых окнах, обучающихся ИИ-систем под дрейфом и цивилизаций с растущим архивом --- обращаться к нестационарному расширению.

\section{Заключение}\label{ux437ux430ux43aux43bux44eux447ux435ux43dux438ux435}

Работа определила витальность операционально как пару \((S, \eta_v)\): предсказательную эффективность самомоделирования \(\eta_v = I_{\text{pred}} / I_{\text{mem}} \in [0,1]\) --- долю удерживаемой системой памяти о среде, которая ещё предсказывает её будущее, --- \textbf{при} замкнутой петле самооплаты \(S\) (§ 2.7). Граница \(\eta_v \le 1\) для пассивного самомоделирования --- следствие неравенства об обработке данных (без отдельной термодинамической посылки; активный случай --- § 2.1), не энергетическая нормировка; термодинамическую цену удержания устаревающей памяти несёт отдельный якорь --- бонд Still (2). Структурный скрининг требует одновременного присутствия трёх ёмкостей --- \(T_v\), \(C_v\), \(S_v\), --- собранных в композитный индекс структурной готовности \(I_v\). Индекс --- необходимое, но не достаточное условие витальности; он занимает то же позиционное место, что assembly index, \(\Phi\), морфодинамика Дикона. Сама \(\eta_v\) --- операциональный критерий синхронной витальности в смысле бинарного знака (\(\eta_v > 0\) при замкнутой петле самооплаты). Критерий отличает структурно богатые, но невитальные системы (Солнце, ураган, LLM при границе «веса в момент инференса» --- \(\eta_v^{\text{int}}\) не определён) от систем со структурно замкнутой петлёй самооплаты по counterfactual ablation § 2.2. Среди систем со замкнутой петлёй бактерия имеет воспроизводимо измеренный \(\eta_v > 0\) через § 3.5. Город, биосфера, LLM-как-корпорация --- со структурно положительным знаком (\(S = 1\)) при количественно не определённой величине \(\eta_v\) (MI-оценка \(I_{\text{pred}}\) и \(I_{\text{mem}}\) --- открытая задача, § 3.4, § 3.6, § 3.7). Содержателен межклассово именно знак, не абсолютная величина (§ 3.5).

Применение двухступенчатой процедуры к шести камертонам даёт резкие демаркационные ответы и делает явной структурную асимметрию: сложность повсеместна, замкнутые петли самооплаты редки. Сопоставление с assembly theory, IIT, телеодинамикой и FEP позиционирует \(\eta_v\) как операциональное дополнение каждой из программ. Сформулированы конкретные условия фальсификации, привязанные к проверяемым физическим, информационным и эволюционным предсказаниям.

Программа открыта к опровержению по нескольким независимым линиям, с эмпирическим риском, концентрированным в ложно-отрицательном проникновении и несущем Прогнозе 1 (§ 5; § 2.7, асимметрия фальсификации; с отдельно фиксируемым условием подтверждения уникальности операционализации), и дисциплинирована против спасательных манёвров требованием называть конкретную провалившуюся ёмкость при каждом обнулении (§ 2.6): опровержение одной ёмкости обрушивает обе шкалы. Эта дисциплина ограничивает, но не устраняет лакатошевский риск защитного пояса (§ 2.6, § 2.7). В этом её научный статус --- не претензия на окончательное определение жизни, а проверяемая исследовательская программа с явно артикулированными границами применимости.

\emph{Биологические следствия.} Главные содержательные эффекты \(\eta_v\)-программы --- внутри теоретической биологии, не на астробиологической или AI-периферии. (а) Синхронная/диахронная декомпозиция витальности (§ 2.7) даёт новую формулировку вопроса «что есть организм» в духе Pradeu (2018) и Bich and Green (2018): синхронный \(\eta_v > 0\) --- необходимое условие индивидуированности живого, диахроническая дарвиновская динамика --- её дополнительная размерность. (б) Минимальные клетки (JCVI-syn3.0, Hutchison et al.~2016) и вирусы получают операционализируемую границу через counterfactual ablation § 2.2: пара (WT \emph{E. coli} / syn3.0) различима не качественным критерием organisational closure, а количественной шкалой \(\eta_v\) через оба компонента (§ 4.7, § S4.7), а вирусные частицы попадают под переадресацию субъекта на (virion + host) --- диагноз структурно гомологичный LLM-corp (\(B_4\), § 3.4). (в) Major transitions in evolution (§ 4.8, § S4.8 --- после гиперцикла и major transitions Maynard Smith--Szathmáry) переформулируются как события переадресации субъекта self-payment с термодинамическим критерием детектирования события. Появляется новый агрегатный актор с собственным \(I_{\text{pred}}^{\text{new}}\), оплачиваемым диссипацией на новом уровне. (г) Симбиотические системы --- растение--микориза (Simard et al.~1997; ср. критическую переоценку Karst et al.~2023), коралл--zooxanthellae --- задают испытательный стенд для дисциплины конвенциональности границы § 2.2: чей \(\eta_v\) растёт при симбиозе, как зависит знак от выбора границы между партнёрами, при каких условиях симбиотический контур формирует новый субъект self-payment. (д) Биоэлектрический морфогенез (Levin 2019) --- биоэлектрические паттерны как \(M_X\) (физическая реализация модели морфологии), оплачиваемые ионными токами через клеточные мембраны, --- даёт прямой кандидат на \(\eta_v\)-измерение в морфогенетическом поле. Степенная зависимость темпа эволюционной адаптации от \(\eta_v\) (Прогноз 1, § 5) --- проверяема на одной микробной линии в обычной экспериментальной эволюции, не требуя астробиологического или AI-контекста, и составляет центральное эмпирическое предсказание программы для теоретической биологии.

\emph{Спекулятивные следствия (астробиология, AI-этика, демаркация).} За пределами биологического ядра программа даёт переформулировку традиционно метафизических вопросов. Поиск внеземной жизни получает термодинамическое требование замкнутых петель самооплаты поверх существующих химических биосигнатур (дополнение к Ladder of Life Detection, не замещение). Этический статус искусственных систем ставится как вопрос о появлении собственной диссипативной платы за модель (помимо традиционных линий sentience / welfare / functionalism). Граница между структурно сложным и витальным получает операциональный \emph{необходимый} срез: \(\eta_v > 0\) --- необходимое условие синхронной витальности (§ 2.7); верхняя граница \(\eta_v \le 1\) --- \textbf{безусловная теорема} (следствие неравенства об обработке данных, § 2.1) для пассивного самомоделирования; активный случай (сенсомоторная обратная связь) --- открытая задача (направленная информация, бонд Sagawa--Ueda 2010). Полное определение жизни в смысле NASA-формулы (Joyce 1994) требует также дарвиновской динамики на популяционном уровне, и обе размерности --- синхронная и диахронная --- взаимодополнительны. Каждая из этих переформулировок не есть автоматическое решение соответствующего вопроса; она есть указание на конкретную измеримую величину, относительно которой вопрос можно формулировать дисциплинированно.

Для теоретической биологии \(\eta_v\) открывает количественный мост между традицией organisational closure (autopoiesis, M,R-системы) и предсказательными моделями адаптации.


\backmatter

\section*{Декларации}

\bmhead{Финансирование}
Исследование не получало внешнего финансирования.

\bmhead{Конкурирующие интересы}
Автор заявляет об отсутствии каких-либо релевантных финансовых или нефинансовых интересов, которые могли бы быть восприняты как влияющие на содержание настоящей статьи (за период последних трёх лет до подачи рукописи).

\bmhead{Этические декларации}
Этическое одобрение, согласие на участие и согласие на публикацию неприменимы: настоящая работа --- теоретическое исследование с численными симуляциями, не вовлекавшее людей, животных, биологический материал или связанные с ними данные в качестве субъектов.

\bmhead{Доступность данных}
Первичные эмпирические данные для этой работы не генерировались. Воспроизводимый код разобранных примеров (§§ 3.4--3.7, 4.4: \emph{E. coli}-хемотаксис, биметаллический термостат, синтетические городские профили, биосфера-как-компрессор и LLM-как-корпорация) открыто доступен в GitHub-репозитории \texttt{andriishin/landauer-self-modeling-oa} (https://github.com/andriishin/landauer-self-modeling-oa) и вместе с препринтом этой рукописи заархивирован на Zenodo (DOI: 10.5281/zenodo.21039160).

\bmhead{Вклад авторов}
Концептуализация, методология, формальный анализ, написание --- подготовка первоначального черновика, написание --- рецензирование и редактирование: А.А. Автор прочёл и согласился с опубликованной версией рукописи.

\bmhead{Использование инструментов ИИ}
В соответствии с Springer Nature editorial policy on AI tools in scholarly writing ИИ не указывается как соавтор и автор сохраняет полную ответственность за всё содержание рукописи. При подготовке настоящей рукописи и сопроводительных материалов автор использовал Anthropic Claude (Opus 4.6, 4.7, 4.8) для пяти задач. (i) Генерация и верификация воспроизводимого Python-кода разобранных примеров (§§ 3.4--3.7, 4.4; код --- в репозитории данных, см. раздел «Доступность данных»). (ii) Сверка библиографических метаданных с локальными копиями цитируемых источников. (iii) Редакторская правка человеко-порождённого текста на уровне читабельности, грамматики и стиля (copy-editing); согласно руководству Springer Nature, такое использование не требует отдельной декларации, но указывается здесь для полной прозрачности. (iv) Перевод рукописи и сопроводительных материалов (включая сопроводительный README кода) с русского первоисточника на английский --- язык подачи журнала --- с последующей терминологической и стилистической ревизией; автор проверил перевод и несёт за него полную ответственность. (v) Помощь в выводе и верификации формального аппарата (вывод границы \(\eta_v \le 1\) как следствия DPI и термодинамического бонда Still (2)) под руководством автора; автор независимо проверил все доказательства и несёт за них полную ответственность. Все научные утверждения --- формулировка \(\eta_v\), требование self-payment, диагнозы paradigm-cases и критерии фальсификации --- замышлены автором; их формальное оформление выполнено с ИИ-ассистированием под контролем автора. ИИ-инструменты не порождали научного содержания автономно. Автор проверил каждый ИИ-ассистированный фрагмент и принимает полную ответственность за содержание рукописи и сопроводительных материалов.


\nocite{*}
\bibliography{refs}

\end{document}
